\documentclass[11pt]{article}

% ----------- Packages -----------
\usepackage[margin=1in]{geometry}
\usepackage{setspace}
\usepackage{amsmath,amssymb,amsthm}
\usepackage{booktabs}
\usepackage{multirow}
\usepackage{graphicx}
\usepackage{caption}
\usepackage{subcaption}
\usepackage{natbib}
\usepackage{hyperref}
\usepackage{xcolor}
\hypersetup{colorlinks=true, linkcolor=blue!50!black, citecolor=blue!50!black,
  urlcolor=blue!50!black}
\usepackage{footmisc}
\usepackage{enumitem}
\usepackage{titlesec}

\titleformat{\section}{\large\bfseries}{\thesection.}{1em}{}
\titleformat{\subsection}{\normalsize\bfseries}{\thesubsection.}{1em}{}

\onehalfspacing
\setlength{\parindent}{2em}

\newtheorem{proposition}{Proposition}
\newtheorem{definition}{Definition}

\title{\Large \textbf{A Computational Laboratory for Industry Dynamics:\\
Multi-Agent Large-Language-Model Simulation of an Oligopolistic Biotech Industry}}
\author{\normalsize Anonymous\thanks{We thank seminar participants at \textit{(omitted for review)}.
All errors are our own. Replication code, prompts, and the simulation environment
will be made publicly available upon acceptance.}}
\date{\normalsize \today}

\begin{document}
\maketitle

\begin{abstract}
\noindent We introduce a computational laboratory in which the strategic
decisions of firms, financial intermediaries, regulators, executives, and
activist investors are produced by separately-instantiated large language
model (LLM) agents whose information sets reflect the same disclosure
boundaries that organize real economies. Firms decide on production,
pricing, R\&D, marketing, capital structure, and dividends; commercial banks
underwrite revolvers and term loans; investment banks execute equity and
bond issuances; an equity market priced by a panel of LLM analysts marks
shares each quarter; an environment agent resolves the product market;
boards governed by hidden CEO types hire and fire executives; activists
launch campaigns; and a private-equity sponsor funds entry. We report
results from a 20-year (80-quarter) simulation of an emerging longevity-drug
industry seeded with six firms and exposed to free entry. Seventeen distinct
firms compete over the run; eight exit through bankruptcy and three are
absorbed via distressed auction. The Herfindahl index averages 1{,}697,
the average top-firm share is 24.2\%, and the cumulative failure rate
reaches 47\%. We document six economic findings that the laboratory delivers from
emergent agent reasoning rather than from programmed update rules.
(i) Firm mortality is bimodal — a leverage-trap cluster of overconfident
expansion failures in early life, a slow-bleed cluster of position-erosion
failures in mid-life — and the bimodality is identifiable in the firms'
strategic memos in a way conventional structural models cannot resolve.
(ii) A roll-up acquirer emerges endogenously: a single solvent firm
pre-funds itself for peer distress, deploys cash at three successive
auctions, and absorbs the residual assets of failed competitors at
fire-sale prices, consistent with \citet{shleifervishny1992}.
(iii) Activist buyback campaigns systematically fail because cash is a
recurrent real option in the \citet{dixitpindyck1994} sense and the
activist's signal is a single-shot demand. (iv) LLM CFOs systematically
under-lever relative to the trade-off optimum because the rhetoric of
covenant downside is more available in the cognitive substrate than
the rhetoric of marginal tax-shield benefit — a testable prediction in
which augmenting the prompt with a worked tax-shield example should
shift behaviour. (v) The PE channel functions as Schumpeterian
renewal: it is the financial intermediary's decision rule that prevents
the industry from collapsing to monopoly, with the rate of competitive
turnover monotonic in PE underwriting capacity. (vi) Equity-market
reflexivity emerges from the panel's correlated anchoring tendency
combined with the firm's narrative interpretation of price moves,
generating a Soros-style price--fundamental feedback loop without
either side being programmed to produce it. We discuss what the
laboratory does and does not deliver, the engineering choices that
determine which findings should be taken as substantive versus
artifactual, and the methodological position LLM industrial
organization occupies between agent-based modeling, structural
estimation, and behavioural theory of the firm.
\bigskip

\noindent\textit{JEL Codes:} L11, L13, G33, C63, D83.
\smallskip

\noindent\textit{Keywords:} industrial organization, agent-based models,
large language models, simulation, market structure, financial distress.
\end{abstract}

\thispagestyle{empty}
\newpage

% ============== 1. INTRODUCTION ==============
\section{Introduction}

Industrial organization has long alternated between two unsatisfying
empirical positions. Reduced-form analysis of historical industry data
delivers credibly identified estimates of narrow effects but cannot
generate counterfactual industries; structural estimation can simulate
counterfactuals but only after one has imposed the strong functional-form
and equilibrium-selection restrictions that make a likelihood
tractable.\footnote{See \citet{ackerberg2007econometric}, \citet{pakes2003}, and
\citet{aguirregabiria2010dynamic} for canonical statements of this trade-off.}
Game-theoretic models of strategic firm behavior, in turn, populate worlds
with optimizers whose decision rules bear little resemblance to the
heuristic, narrative-driven, constraint-saturated decisions executives
actually make. The Nelson--Winter \citeyearpar{nelsonwinter1982} evolutionary
program responded to this disconnect by replacing maximization with
routines, and the agent-based literature in the tradition of
\citet{tesfatsion2006agent}, \citet{axtell2001zipf}, and
\citet{farmer2009economy} extended the response with computational
heterogeneity. Yet evolutionary and agent-based firms remain stylized: a
small number of state variables, hand-crafted update rules, and---most
limiting---no language. Real firms communicate, justify, deflect, and
persuade. Boards interrogate strategy. Investors react to disclosure. None
of that is in the agent-based menu.

This paper proposes an alternative. We assemble a multi-agent simulation
in which the agents that drive the industry are large language models
(LLMs) instantiated separately for each role, each prompted with its own
information set and instructions, and each issuing structured decisions
that the simulation engine commits to a balance-sheet-consistent state.
The laboratory contains, at minimum, eight kinds of agents: (i) operating
firms, who decide production, pricing, capital expenditure, R\&D,
marketing, financing, and payout policy; (ii) an industry environment, a
``referee'' that resolves the product market, grants R\&D outcomes, and
narrates idiosyncratic shocks; (iii) a commercial bank that underwrites
revolvers and amortizing term loans subject to covenants; (iv) an
investment bank that prices and places equity issuances and bond
offerings; (v) an equity market priced by a panel of LLM valuation
analysts whose median quote becomes the marked share price; (vi) a
private-equity sponsor that funds the activation of dormant entrants;
(vii) boards of directors with discretion to dismiss the CEO, configured
so that the CEO has a hidden behavioral type unobserved by the board;
(viii) activist investors who occasionally launch campaigns---typically
buyback demands aimed at cash-rich firms. Optional modules implement
auditors, sell-side analysts, an SEC enforcement agent, M\&A bidding,
restatements, and an earnings-management injection process; we run with
the production-and-finance subset active.

The combination is, on its face, alarming. LLMs hallucinate; they are
sensitive to prompt formulation; they have no calibration to first
principles. We share these concerns and address them throughout the paper:
in the architectural separation of decisions from accounting, in the
balance-sheet invariants that fail loudly when violated, in the
deterministic clamps that prevent any agent from generating physically
infeasible outcomes, and in the panel-median valuation procedure designed
to suppress single-LLM outliers. The laboratory should be understood not
as a forecasting tool but as a controlled environment for studying which
properties of industry dynamics are emergent from a particular cognitive
architecture, in the sense that the regularity does not have to be
written into any agent's prompt to appear in the data the simulation
generates.

A concrete example may help fix ideas. In our 80-quarter run, firm\_14
spawns at quarter 46 with a high-baseline capability stock,
representing a science-led entrant backed by a private-equity sponsor.
Within three quarters of activation it has moved into the industry
lead at a 22\% market share, displacing the incumbent firm\_7. The
incumbent's CFO, in his Q49 strategic memo, writes: \emph{``A
new entrant (firm\_14) has emerged with what appears to be a credible
clinical-stage advantage. We are not in a position to match their
capability stock through R\&D alone within a reasonable horizon. Our
response strategy is to (i) protect our installed base through
aggressive marketing, (ii) preserve our cash position for the
distressed-asset opportunities that the next two-to-three years are
likely to produce as marginal competitors face dual pressure from
firm\_14 and from rising R\&D costs, (iii) defer Generation-2
investment until firm\_14's leapfrog has been validated or
falsified.''} The memo prefigures, in narrative form, the M\&A
consolidation strategy that firm\_7 actually executes over the
subsequent twenty-five quarters: firm\_7 acquires three distressed
peers between Q44 and Q71. The memo is not a coefficient. It is the
firm's articulated rationale for a sequence of decisions that
unfolds over six years, and the simulation records both the rationale
and the decisions in machine-readable form. This is the kind of
output that, to our knowledge, no other methodology in industrial
organization is capable of producing.

We make three contributions. First, \textit{economically}, we present
a set of findings about industrial dynamics that the laboratory
generates emergent and that connect to the standard theoretical
literature. Mortality is bimodal: an early-life cluster matches the
\citet{jensen1986agency} free-cash-flow over-investment story
modulated by debt, while a mid-life cluster reflects gradual
position erosion. Roll-up consolidation through \citet{shleifervishny1992}
fire-sale dynamics emerges from one acquirer's deliberate
pre-funding strategy, articulated in advance in its strategic memos.
Cash hoarding by survivors is rationalised in
\citet{dixitpindyck1994} real-options language and resists activist
pressure because the activist's demand is one-shot while the option
value is recurrent. LLM CFOs systematically under-lever because
covenant rhetoric is cognitively more available than tax-shield
rhetoric, generating a testable prediction about the response of
firm leverage to prompt-augmented financial reasoning. The PE
channel operates as the Schumpeterian renewal mechanism that keeps
the industry from collapsing to monopoly. CEO turnover by hidden
agency-cost type tracks \citet{jensen1986agency} predictions about
observable behavioural signatures. Each of these findings has the
property that it would be hard to obtain in conventional
agent-based simulation (which lacks narrative content) or in
structural estimation (which assumes the optimization
characterisation we are trying to test against).

Second, \textit{methodologically}, we describe an
information-theoretic separation between agents---each prompt builder
takes only the data the corresponding real-world actor would have---that
appears to be the binding constraint on whether the simulation produces
recognizable industrial dynamics rather than a coordinated narrative
hallucination. We document concrete failure modes (e.g., dataclass-default
zero contamination, environment over-anchoring, single-LLM equity
outliers) and the engineering responses that brought them under control.
Third, \textit{conceptually}, we situate the LLM laboratory between
agent-based modeling, structural IO, and the behavioral theory of the
firm, arguing that its comparative advantage is the simulation of
\emph{decision processes} rather than \emph{decision rules}. The
substantive consequence is that the laboratory can produce
\emph{mechanism evidence}: the strategic-memo data identifies
\emph{which} of several candidate economic mechanisms is operating in
any given firm-quarter, in a way that conventional industrial-organization
data cannot. Standard reduced-form analysis observes outcomes; the
laboratory observes outcomes alongside the firm's articulated
reasoning for the decision that produced them.

A reasonable response, at this point, is that we are reinventing
agent-based macroeconomics with more expensive parts. We disagree, and
the disagreement structures the paper. Agent-based models typically
specify a parametric decision rule and analyze the resulting dynamics;
LLM-driven models do something materially different, namely produce a
\emph{decision process} from textual reasoning over a structured
information set. The distinction is not cosmetic. A parametric rule
inverts back to the parameters; a decision process inverts back to
narrative content---to phrases like ``we believe Generation 2 readiness
remains 18--24 months out and have therefore deferred capacity
expansion.'' That phrase is not a coefficient. It is a
machine-interpretable statement about an executive's reasoning, and it
can be analyzed by the same machinery that produced it. We will return
to this point in Section \ref{sec:mechanisms}; for now we note only that
the laboratory's auditable narrative output is its principal
methodological surplus over the agent-based and structural alternatives.

A second reasonable concern is reproducibility. LLMs are stochastic, and
LLM providers update model weights without disclosing the change. We
address this through three layers of discipline. First, the orchestrator
is deterministic given a random seed and a model roster, and the
roster pins specific model versions; the same roster on the same seed
will reproduce within the model provider's own determinism budget.
Second, every LLM call is logged with its exact prompt and response,
producing a per-quarter \texttt{llm\_calls.jsonl} that is sufficient to
reconstruct the deterministic state evolution after the fact, even if
the underlying model weights have changed. Third, every quarter
boundary is a snapshot, and any conclusion of this paper can be
recomputed from the snapshot directory without re-querying any LLM. The
combination is reproducibility-with-an-asterisk, weaker than that
provided by closed-form models but considerably stronger than that
provided by uncached LLM applications.

The remainder of the paper proceeds as follows. Section \ref{sec:literature}
positions the paper relative to four literatures: agent-based computational
economics, structural industrial organization, the behavioral theory of
the firm, and the emerging literature on LLM agents in economic
applications. Section \ref{sec:environment} describes the simulation
environment in technical detail, including the agent menu, the information
boundaries, the formal demand and cost model, the accounting and
balance-sheet invariants, the product market, and the financing and
bankruptcy mechanics. Section \ref{sec:design}
describes the experimental design: model roster, calibration, parameter
choices, and run protocol. Section \ref{sec:measurement} describes the
output panels and how key quantities are constructed. Section
\ref{sec:results} presents the main results: industry life-cycle,
mortality hazards, concentration dynamics, M\&A, leverage, governance,
equity-market behavior, cash hoarding, activist campaigns, growth
distributions, and detailed firm case studies. Section \ref{sec:mechanisms}
isolates the mechanisms underlying the headline findings and discusses
threats to external validity. Section \ref{sec:conclusion} concludes.
Appendix \ref{app:prompts} reproduces the most consequential prompts.

\section{Related Literature}\label{sec:literature}

\paragraph{Agent-based computational economics.} The closest methodological
relatives of our laboratory are the agent-based models surveyed by
\citet{tesfatsion2006agent}, \citet{leijonhufvud2006agent}, and
\citet{epstein2006generative}. The Santa-Fe-style stock market
(\citealp{lebaron1999time}, \citealp{arthur1997asset}) and the K\&S model of
\citet{dosi2010schumpeter} are the canonical examples within macro and IO,
respectively. The fundamental departure of our work is that decision rules
in those models are hand-crafted update equations parameterized by a
small number of behavioral coefficients. In ours, decision rules are the
emergent output of language models prompted with information packages
that approximate what the corresponding real-world actor would observe.
This trades parametric transparency for a different kind of fidelity: an
LLM-driven CFO that carries a heavy debt load can argue, in writing, why
deferring R\&D is the prudent action this quarter, and the simulation
records the argument; one cannot ask a coefficient why it took the value
it took.

\paragraph{Structural industrial organization.} Our work is complementary
to the dynamic-oligopoly literature initiated by \citet{ericsonpakes1995},
extended by \citet{bajaribenkardlevin2007} and \citet{doraszelskipakes2007},
and surveyed by \citet{aguirregabiria2010dynamic}. Structural IO recovers
deep parameters from observed equilibrium play; we generate equilibrium
play given a particular cognitive architecture, and the
``deep parameters'' of our model are the prompts. In a structural model
the analyst cannot ask whether firms would behave differently if they
reasoned about debt covenants narratively rather than algebraically; in
our laboratory, this is essentially the only kind of question one can
ask. We view the two methodologies as complements: structural estimation
disciplines theory by data, computational laboratories discipline data by
counterfactual.

\paragraph{Behavioral and evolutionary theory of the firm.} The
\citet{cyertmarch1963} behavioral theory and the \citet{nelsonwinter1982}
evolutionary tradition argue that the firm is best modeled as a
collection of routines, aspiration levels, and political coalitions
rather than as a profit maximizer. Our LLM firms inherit something of
this spirit: their decisions are bounded by reasoning capacity, anchored
on previous decisions, and influenced by the framing of their information
package. Where Cyert and March wrote prose about how managers \emph{think},
we now have a tractable substrate that allows us to study what such
thinking generates at scale.

\paragraph{LLMs in economic applications.} A small but growing literature
uses LLMs as economic agents. \citet{horton2023large} proposes
``homo silicus'' as a substitute for human subjects in lab experiments;
\citet{argyle2023out} document that LLMs reproduce demographic
correlations in survey responses; \citet{brand2023using} apply LLMs to
preference elicitation. Closer to our setting,
\citet{ross2024llm} use LLMs as principal-agent decision-makers in
contracting environments, and \citet{li2023}\nocite{li2023}
implement LLM-based macro households. To our knowledge, ours is the first
attempt to assemble a fully populated industry---firms, financial
intermediaries, regulators, governance, capital markets---out of LLM
agents with isolated information sets, and to use the resulting
environment to study industrial dynamics over a 20-year horizon.

\paragraph{Financial frictions and corporate dynamics.} The role of
financial intermediaries in our laboratory connects to the literature on
financial frictions (\citealp{bernankegertler1989}, \citealp{kiyotakimoore1997},
\citealp{rampinjviswanathan2010}), the costly external finance literature
(\citealp{myersmajluf1984}), and dynamic capital-structure models
(\citealp{lelandtoft1996}, \citealp{strebulaev2007}). The capital-structure
results we report are not estimated; they are emergent from how LLM
firms reason about debt under our prompt structure. They should be read
as descriptive facts about that cognitive architecture rather than as
estimates of structural parameters.

\paragraph{Industry dynamics and firm size distributions.} Our findings
on mortality, concentration, and growth distributions speak to a long
tradition. \citet{gibrat1931} conjectured that growth rates are
independent of size; \citet{hall1987relationship} and
\citet{evans1987tests} document violations of Gibrat's law in cohort
studies of U.S. manufacturers; \citet{sutton1997gibrat} surveys the
empirical regularities. Our laboratory's growth distribution is
right-skewed and fat-tailed, with a median quarterly growth rate of
zero and a long upper tail driven by entrants ramping production. We
report the moments and discuss whether they reflect endogenous
dynamics or are artifacts of the entry calibration in
Section~\ref{sec:results}. \citet{caves1998industrial} and
\citet{geroski1995entries} catalogue the entry-and-exit literature; we
find a hazard pattern broadly consistent with theirs but front-loaded
relative to U.S. data, which we attribute to the laboratory's
compressed time scale.

\paragraph{Bankruptcy and distressed asset markets.} The distressed
auctions in our laboratory are mechanically similar to the
\citet{shleifervishny1992} fire-sale mechanism, with a key
difference: prices are set by an LLM auctioneer rather than by an
explicit demand curve. \citet{strömberg2000} documents the
characteristics of Chapter 7 versus Chapter 11 outcomes; the present
laboratory implements only a Chapter-7-like asset auction and we
discuss the resulting bias toward overstated liquidation outcomes. The
\citet{pulvino1998effects} cyclical patterns in distressed prices
emerge in our laboratory: the three M\&A transactions occur during
periods of widespread industry distress, and the prices reflect the
acquirer's strategic-fit narrative more than a clean discount on book
value.

\paragraph{Executive turnover and governance.} The CEO-turnover patterns
we observe connect to \citet{warner1988stock},
\citet{denis1997ownership}, \citet{taylor2010ceo}, and the
\citet{kaplanminton2012} update on hazard rates. Our turnover frequency
is higher than the empirical benchmark, which we attribute to the
laboratory's annual board-review cadence and to the absence of any
patience parameter in the board prompt; tuning these is an obvious
next step.

\paragraph{Computational laboratories more broadly.} The use of
computer simulation as a vehicle for inference about social systems
has a substantial intellectual history. \citet{shubik1964} and the
Cowles Commission tradition exhibit early models of strategic
interaction through computation. The \citet{schelling1971dynamic}
segregation model is the canonical demonstration that emergent
macroscopic order can arise from microscopic individual decisions
that contain no preference for the macroscopic property in question.
Our laboratory is, in a sense, a Schelling-style demonstration with
language: industrial concentration emerges from no firm prompted to
seek it, financial distress emerges from no agent instructed to
generate it, and the regularity these emergent patterns exhibit is
the principal evidence we offer that the laboratory is producing
genuine economic dynamics rather than coordinated fiction.

\section{The Simulation Environment}\label{sec:environment}

\subsection{Architectural overview}

The laboratory is organized around a single mutable object,
\textit{state}, that is the consensus record of who has decided what,
which contracts are outstanding, and how every dollar of every firm's
balance sheet is accounted for. The state is advanced one quarter at a
time by an \textit{orchestrator} that calls each agent in a fixed
sequence, gathers structured outputs (typically JSON with a documented
schema), validates them against deterministic constraints, and commits
the resulting deltas back to the state. Crucially, no agent ever writes
directly to state. Every state mutation is performed by accounting code
that takes the agent's intentions as inputs and enforces conservation of
balance-sheet identities as outputs.

The orchestrator is deterministic given a random seed, the model roster,
and any cached LLM responses. A snapshot of the entire state is written
to disk after each quarter, so a run can be resumed from any quarter
boundary.\footnote{The snapshot mechanism was important in our development.
A 20-year run takes roughly 21 hours of wall-clock time on a workstation
calling the OpenAI and Anthropic APIs; the ability to resume from a
mid-run snapshot after a code fix saved many days of recomputation.}

Formally, denote the state at the close of quarter $t$ by
$\Omega_t = \{\Omega_t^F, \Omega_t^C, \Omega_t^M, \Omega_t^I\}$, where
$\Omega_t^F$ collects the firm-level state (capacity, capital stocks,
balance sheet, executive identity), $\Omega_t^C$ collects the contract
state (debt facilities, equity issuances, M\&A bids), $\Omega_t^M$
collects the macro state (policy rate, market risk premium, political
uncertainty), and $\Omega_t^I$ collects the information state (analyst
notes, audit opinions, activist campaigns). The transition from
$\Omega_t$ to $\Omega_{t+1}$ is the composition of a sequence of
agent-resolution maps,
\[
\Omega_{t+1} \;=\; \Phi_{\text{record}} \circ
\Phi_{\text{settle}} \circ \Phi_{\text{eqmarket}} \circ
\Phi_{\text{ibank}} \circ \Phi_{\text{cbank}} \circ
\Phi_{\text{accounting}} \circ \Phi_{\text{environment}} \circ
\Phi_{\text{firm}} \circ \Phi_{\text{macro}} \;(\Omega_t),
\]
where each $\Phi_{x}$ is a map that takes the state, calls a set of LLM
agents conditioned on a slice of state appropriate to that role, parses
the structured outputs, and returns an updated state. The composition
is deterministic in the sense that, fixing a seed and a cache, the
output state is reproducible. The composition is interpretable in the
sense that each $\Phi_x$ can be inspected, replayed, or replaced.

\subsection{The firm agent}\label{sec:firmagent}

A firm is, internally, a frozen dataclass whose fields are the canonical
operating, financing, and intangible state variables: capacity (in
units), gross PPE and accumulated depreciation, capability and brand
stocks, cash, long-term debt, revolver balance, inventory at FIFO cost,
retained earnings, and so on. These are advanced through a
\texttt{.evolve()} method that takes a dictionary of named field updates
and returns a new firm; mutation is never destructive. This convention
means that any agent's decision either commits cleanly to the next-period
state or fails the accounting check, in which case the firm is rolled
back to a deterministic fallback consistent with the previous quarter.

Each quarter, the firm agent receives an information package containing
(i) a snapshot of its own internal state, (ii) the four-quarter rolling
public Compustat panel for all firms in the industry (revenue, cost of
goods sold, gross profit, operating income, net income, cash, total
assets, leverage, EPS, share price, market cap), (iii) a strategic memo
from a board discussion held at the start of the quarter, (iv) a
covenant compliance summary, (v) an internal credit-officer briefing if
a debt facility is outstanding. It returns a JSON document specifying
production volume, pricing, R\&D and marketing budgets, capital
expenditure, dividend per share, repurchase target, debt issuance or
amortization, equity issuance request, and a free-text strategic
narrative that contextualizes the quantitative choices.\footnote{The
narrative is not just decoration. It is read into the next quarter's
information package as the firm's own ``prior memo,'' giving the
LLM-driven firm a working memory that survives across calls. The
narrative is also the input to the board's CEO performance review.}

We add three structural prompts that bear on the present results.
First, we instruct firms to debate, when their cash position is
substantial, three options---hold for strategic optionality, deploy into
the business, or return to shareholders---and to articulate which they
choose and why. We do not impose a numerical threshold for what counts as
``substantial''; the firm makes the call. Second, we ask firms to
maintain a one-quarter and one-year EPS guidance number that is then
visible to analysts and the equity market. Third, we instruct firms that
mandatory cash obligations (interest, principal due, lease payments)
take precedence over any other use of cash; failure to meet them
triggers the bankruptcy procedure described in
Section~\ref{sec:bankruptcy}.

\subsection{The environment agent}\label{sec:envagent}

The environment agent plays the role of an industry referee. Each
quarter, after firms have submitted their decisions and after a
deterministic clamp has enforced physical feasibility (capacity bounds,
non-negativity, etc.), the environment is shown all firm decisions and
the underlying state, and is asked to produce a verifier output that
specifies, for each firm: realized unit demand, realized R\&D
outcomes (capability gain, optional generation advance, optional
delivery-system advance, optional process-cost reduction), realized
brand-stock change, and any narrative event (regulatory action, supplier
disruption, lawsuit, idiosyncratic productivity shock).

Two design choices about the environment matter empirically. First, the
environment sees the entire industry; it has god-mode visibility. This is
deliberate. A real industry equilibrium is the joint outcome of correlated
draws on capability, brand, demand, and luck; producing them
independently from per-firm prompts gives a different distribution. By
asking one agent to assign all outcomes simultaneously, we get
cross-sectional realism (firms cannot all have lucky quarters at the
same time, market shares add to one, total demand is mediated by a single
market-wide intercept) at the cost of introducing the environment as a
single point of failure if it hallucinates. We address that risk through
deterministic clamps and through a balance-sheet invariant check that
catches any environment output inconsistent with the previous-quarter
state.

Second, R\&D outcomes are returned as guidance rather than fiat. The
environment is told that the cumulative R\&D spend at which a generation
advance becomes plausible is in the neighborhood of \$200 million per
firm but that this number is approximate---firms reach the next
generation at different cumulative levels depending on the quality of
their R\&D, the talent on their team, and regulatory luck. The
environment is instructed to spread advances over time and to narrate the
specific catalyst (Phase 3 readout, regulatory approval,
lead-compound milestone). A by-product of this instruction, which we
discuss in Section~\ref{sec:results}, is that across our 80-quarter run
no firm achieves a generation advance, suggesting that the environment
remains conservative even when several firms have crossed the
\$300--500 million cumulative R\&D mark.

\subsection{Financial intermediaries}

Three intermediaries operate in parallel: the commercial bank, the
investment bank, and the equity market.

\paragraph{The commercial bank.} The commercial bank receives, each
quarter, a request for revolver capacity or an amortizing term loan from
any firm that has decided to draw or to negotiate. It sees the firm's
public Compustat panel, an internal credit file (private credit history
with this lender, covenant compliance trajectory, recent waivers), and
its own portfolio of outstanding facilities. It produces a credit
decision: approve or decline; if approve, principal, term, all-in spread,
covenant package (typically interest coverage and total-debt-to-EBITDA),
and any conditions precedent. Drawn balances accrete interest at the
contracted spread over a deterministic policy rate; covenant violations
are evaluated each quarter and may trigger a revolver pull.

\paragraph{The investment bank.} The investment bank prices and places
equity issuances (rights offerings or shelf takedowns) and bond
issuances. For equity, the inputs are the issuance request, recent
trading prices, and the analyst consensus if sell-side coverage is
enabled. For bonds, the inputs are the firm's leverage, interest
coverage, and a fictionalized credit rating produced jointly with the
issuance decision. The bank produces a placement spread (the
underwriting discount) and an indicative coupon.

\paragraph{The equity market.} The equity market produces a single mark
each quarter: the price at which the firm's shares are deemed to trade.
Producing this price reliably turned out to be the single most
engineering-intensive part of the simulation. With a single LLM, we
encountered occasional unphysical quotes (a quintupling of the share
price on declining fundamentals, a dollar-cents typo, a misread of the
prior-quarter price). We solved this by replacing the single-LLM
valuation with a panel of three valuation analysts---typically the
LLM cores assigned to other firm CFOs, instantiated independently with
the same equity prompt---and taking the per-firm median quote.

The equity prompt provides each panel member with: (i) the four-quarter
rolling price history of every firm in the industry (so the panelist can
read trajectories rather than seeing only a single anchor); (ii) the
firm's most recent management guidance; (iii) the public Compustat panel;
(iv) the analyst consensus if sell-side analysts are enabled. The
panelist returns a target price for each firm, a discount-rate sensitivity,
and a narrative justification. The simulation records all panel votes;
the median is committed to state. The result has been to suppress---though
not eliminate---the most egregious single-LLM excursions.

\subsection{Auxiliary agents}

The remaining agents are conceptually less central but operationally
non-negligible.

The \emph{board of directors} convenes annually (every fourth quarter) and
performs a CEO performance review that has access to a four-year rolling
operating history, the cumulative cost of the CEO's compensation, recent
strategic narratives, and any activist communications. The board can
dismiss the CEO. CEO types (\textit{aggressive\_grower},
\textit{conservative\_steward}, \textit{empire\_builder},
\textit{honest\_operator}) are assigned at firm creation, hidden from the
board, and influence the firm's decision narrative through a system
prompt addendum that the board cannot see.

The \emph{activist agent} surveys all firms each quarter and may launch a
campaign: a buyback demand on a cash-rich firm, a strategic-review
demand on a firm with poor segment performance, or a divestiture demand
on a conglomerated firm. Campaigns enter the firm's information package
the following quarter and its strategic memo for that quarter is required
to address the activist publicly.

The \emph{private-equity sponsor} surveys dormant entrants---firms that
have spawned but not been activated---and decides which to fund. A
PE-funded entrant becomes active with an initial cash injection and a
larger-than-baseline capability stock, simulating a leapfrog entry by an
operator backed by patient capital.

\subsection{Information boundaries}

The most important non-obvious property of the laboratory is that no
agent's prompt builder receives the full state. Each builder takes only
the slice of state the corresponding real-world actor would observe. A
firm sees its own private state plus the public Compustat panel; the
equity market sees the public panel plus analyst notes; the SEC, when
enabled, sees the public panel and any disclosure-triggering events; the
auditor sees its client's books. Manipulation truth (the difference
between true and reported earnings, when earnings management is enabled)
is visible only to the firm that chose it and to the SEC if it elects to
investigate. This information segmentation is what makes the laboratory
behave like an industry rather than like a single agent with a long
context.

\subsection{Accounting and balance-sheet invariants}\label{sec:bsinvariants}

Every decision and every event is intermediated by a small accounting
module that enforces double-entry consistency. Cash, PPE, inventory, and
debt are mutated only through the corresponding journal entries. After
each quarter, an invariant checker recomputes total assets, total
liabilities, and total equity for every firm and verifies that
$A = L + E$ to within a tolerance of \$1{,}000---a rounding band wide
enough to absorb floating-point drift, narrow enough to catch any genuine
omission. Violations are written to a \texttt{bs\_violations.jsonl} log;
during development this log was the principal device for surfacing bugs
that would otherwise have remained silent.

We catalog the four bug classes we encountered, fixed, and pinned with
regression tests. (1) \emph{Dataclass-default zero contamination}: an
exception handler in the firm-decision parallel pool returned a
default-initialized decision object, populating real fields with the
dataclass default of zero. The same six firms hit the exception each
quarter, producing a coordinated zero-production state that initially
looked like a coordination failure; it was, in fact, a software fault.
Fix: explicit carry-forward of the previous quarter's decision when a
fallback is required. (2) \emph{Auction event drop}: the
distressed-auction event-application loop was inside an else-branch
intended for legacy code and so never executed in modern runs.
Three M\&A consolidations had been computed but never committed to
state; the fix outdented the loop. (3) \emph{Auction cash overwrite}: the
defaulted firm's pre-default cash was being overwritten by sale proceeds
rather than augmented by them, silently destroying balance sheet entries.
The fix replaces overwrite with addition and adds a multi-tier creditor
waterfall. (4) \emph{Invariant blind spot}: the balance-sheet check
skipped firms in default, which masked all auction-induced residuals.
Removing the skip produced a noisier log in the short term and a cleaner
post-default panel.

\subsection{The product market}

The product market is differentiated horizontally. Each firm posts a
price and a quantity of units offered for sale; the environment computes
a market-clearing demand for each firm based on its capability stock,
brand stock, list price, and an idiosyncratic match score that varies
across quarters. Capacity binds in both directions: firms cannot sell
more units than capacity provides, and units in excess of demand
accumulate as inventory at FIFO cost. Cost of goods sold is the same
unit cost applied to the same number of units that are recognized as
revenue, ensuring revenue/COGS consistency.\footnote{Earlier versions of
the simulation contained a unit mismatch in which COGS was computed
against a different unit count than revenue; the resulting reported
gross margins were artifactual. Section~\ref{sec:environment} discusses
how this was fixed and pinned.}

The match score is the lever through which the environment
encodes idiosyncratic competitive dynamics: a brand-loyal segment of
the market reaches firms with high brand stock more reliably; a
quality-sensitive segment is drawn to firms with high capability; a
price-sensitive segment chases the lowest list price. The environment
also reasons about regional market saturation and substitution patterns;
the regional dimension is configurable and is on by default in our run.

Although the environment delivers demand allocations as a single
LLM-generated assignment, the underlying mental model the environment
prompt encodes is recognizable. We specify an indicative latent demand
structure of the Hotelling form,
\[
U_{i\theta t} \;=\; \alpha_q \log Q_{it} + \alpha_b \log B_{it}
  - \alpha_p \log p_{it} - \alpha_d \, d(\theta, x_{it}) + \varepsilon_{i\theta t},
\]
where $U_{i\theta t}$ is consumer $\theta$'s utility from firm $i$ in
quarter $t$, $Q_{it}$ is firm $i$'s capability stock, $B_{it}$ is its
brand stock, $p_{it}$ is its price, $x_{it}$ is the firm's position in
a horizontal product space, $\theta$ is the consumer's location in
that space, and $\varepsilon_{i\theta t}$ is an idiosyncratic Gumbel
shock. Quantity demanded $D_{it}$ is then the integral of the choice
probability over the consumer-mass distribution. We do not impose this
specification; the environment prompt instead instructs the LLM to
think in terms of (i) absolute capability level matters quadratically,
(ii) brand depreciates without continued marketing, (iii) consumers
substitute away from price increases at an elasticity ``somewhere
between $-1$ and $-3$'' depending on category, (iv) regional submarkets
saturate, and (v) idiosyncratic demand shocks are non-zero in
expectation. The environment-produced quantity allocation is then
checked for consistency with the prior quarter's posted prices and
capacity bounds, but is not forced to match the indicative
specification above. We retain the ambiguity deliberately: the
environment is allowed to reflect that latent demand has structural
heterogeneity that is not summarized by a single elasticity.

The cost side is more disciplined. Each firm carries a
\textit{base unit cost} $c_{it}$ that depreciates with capability
through learning by doing: $c_{i,t+1} = c_{it}\,(1 - \rho \, \Delta Q_{it})$
where $\rho$ is a learning rate and $\Delta Q_{it}$ is the
capability increment. Realized cost of goods sold in $t$ is
$c_{it} \times D_{it}$ where $D_{it}$ is the realized number of units
sold. The same $D_{it}$ that drives revenue drives COGS; this matters
because an earlier version of our code carried a unit mismatch
(Section~\ref{sec:bsinvariants}) in which COGS was applied to a
clamped quantity while revenue was applied to the unclamped
declaration. Capacity is bounded above by gross PPE divided by a
unit-capital coefficient, with depreciation at 2.5\% per quarter
applied straight-line. Inventory is carried at FIFO unit cost; unsold
production at the close of $t$ enters inventory at $c_{it}$ and is
released against revenue in subsequent quarters in order of arrival.

The R\&D production function is encoded in the environment prompt
rather than in code. The environment is told that capability gain in
quarter $t$ is increasing in $\mathit{xrdq}_{it}$, the quarterly R\&D spend, with
diminishing returns; that gains are mean-zero across firms (a high
draw for one firm is correlated with a low draw for another, capturing
competitive race dynamics); and that the firm-specific capability
trajectory may exhibit step changes corresponding to generation
transitions. The environment is given the four-quarter trailing R\&D
spend for every firm and is instructed to allocate capability gains in
a way consistent with relative spend and team quality. Generation
transitions are described in the prompt as ``approximately
\$200 million in cumulative R\&D'' with the explicit caveat that this
is guidance and not a hard threshold. Brand-stock evolution follows a
parallel logic: the environment rewards marketing spend ($\mathit{xsgaq}_{it}$)
with a brand-stock gain that depreciates 5\% per quarter without
continued investment.

The market-clearing equation we are implicitly imposing on the
environment is therefore
\[
\sum_{i \in \mathcal{A}_t} D_{it} \;=\; \bar{D}_t(\Omega_t^M, p_{1t}, \ldots, p_{Nt}),
\]
where $\bar{D}_t$ is the aggregate market size, itself increasing in
political certainty (a macro-state lever) and decreasing in the
average price level. The environment is shown $\bar{D}_t$ for the
trailing four quarters and is asked to allocate units across firms
consistent with their relative attractiveness. We have not attempted
to recover $\bar{D}_t$ structurally; doing so would require assuming a
particular utility specification, which would defeat the purpose of
delegating the allocation to a flexible reasoner.

\subsection{Financing and capital structure}

A firm finances operations through retained earnings, drawn revolver,
amortizing term debt, and discretionary equity issuances. The revolver
is governed by a borrowing base (a function of receivables, inventory,
and PPE) that is recomputed each quarter; the term loan amortizes on a
schedule set at issuance. Bond issuances become available once the firm
crosses a notional size threshold and provide a higher leverage ceiling
at the cost of a fixed coupon and a less flexible covenant package.

The investment-bank prompt is asked to consider trade-off and pecking-order
intuitions when pricing securities, but does not impose a target
leverage. We will see in Section~\ref{sec:results} that the equilibrium
leverage in our laboratory is markedly below what trade-off theory would
predict at observed tax rates and bankruptcy costs. This appears to be
emergent behavior of the firm-CFO LLMs, who in their narratives
consistently express a preference for cash buffers over interest-tax-shield
maximization.

\subsection{The orchestrator loop in pseudocode}\label{sec:loop}

To make the architectural separation concrete, we reproduce the
quarter-advance loop in compact pseudocode. The loop is the same for
every quarter; toggles enable or disable specific blocks, but the
underlying ordering is invariant. The pseudocode below is faithful to
the production code modulo cosmetic simplification.

\begin{quote}\small
\texttt{\textbf{def} run\_quarter(state, t, config, agents):}\\
\texttt{~~~~state.macro = $\Phi_\text{macro}$(state, agents.macro)}\\
\texttt{~~~~state = inject\_entry(state, t, config)}\\
\texttt{~~~~\textbf{if} config.ma\_enabled:}\\
\texttt{~~~~~~~~bids = $\Phi_\text{ma\_bid}$(state, agents.ma\_bidders)}\\
\texttt{~~~~~~~~state = apply\_ma\_bids(state, bids)}\\
\texttt{~~~~\textbf{if} config.sec\_enabled:}\\
\texttt{~~~~~~~~state = $\Phi_\text{sec\_pre}$(state, agents.sec)}\\
\texttt{~~~~raw = run\_firm\_decisions\_parallel(state, agents.firms)}\\
\texttt{~~~~clamped = enforce\_physical\_clamps(state, raw)}\\
\texttt{~~~~outcomes = $\Phi_\text{environment}$(state, clamped, agents.env)}\\
\texttt{~~~~state = $\Phi_\text{accounting}$(state, clamped, outcomes)}\\
\texttt{~~~~\textbf{if} config.earnings\_announcement:}\\
\texttt{~~~~~~~~state = $\Phi_\text{announcement}$(state, agents.firms)}\\
\texttt{~~~~\textbf{if} config.sellside\_analysts:}\\
\texttt{~~~~~~~~state = $\Phi_\text{sellside}$(state, agents.analysts)}\\
\texttt{~~~~state = $\Phi_\text{equity}$(state, agents.equity\_panel)}\\
\texttt{~~~~state = $\Phi_\text{ibank}$(state, agents.ibank)}\\
\texttt{~~~~state = $\Phi_\text{cbank}$(state, agents.cbank)}\\
\texttt{~~~~\textbf{if} config.sec\_enabled:}\\
\texttt{~~~~~~~~state = $\Phi_\text{sec\_post}$(state, agents.sec)}\\
\texttt{~~~~state = settle\_obligations\_and\_check\_bankruptcy(state)}\\
\texttt{~~~~check\_balance\_sheet\_invariants(state, t)}\\
\texttt{~~~~\textbf{if} t \% 4 == 0:}\\
\texttt{~~~~~~~~\textbf{if} config.auditor\_enabled:}\\
\texttt{~~~~~~~~~~~~state = $\Phi_\text{audit}$(state, agents.auditors)}\\
\texttt{~~~~~~~~\textbf{if} config.governance\_enabled:}\\
\texttt{~~~~~~~~~~~~state = $\Phi_\text{governance}$(state, agents.boards)}\\
\texttt{~~~~~~~~\textbf{if} config.restatements\_enabled:}\\
\texttt{~~~~~~~~~~~~state = $\Phi_\text{restate}$(state)}\\
\texttt{~~~~snapshot(state, t)}\\
\texttt{~~~~\textbf{return} state}
\end{quote}

The two non-obvious features of this loop are the
\texttt{enforce\_physical\_clamps} step and the
\texttt{check\_balance\_sheet\_invariants} step. The first prevents
any LLM agent from generating a physically infeasible decision (a
negative production volume, a price below zero, capital expenditure
in excess of available cash plus revolver capacity); when a clamp
binds, the resulting clamped value is propagated to all downstream
accounting. The second is the loud-failure principle in operation:
any inconsistency between assets and liabilities-plus-equity, beyond
a tolerance of \$1{,}000, is logged with the offending firm-quarter
and a snapshot of the state at the time of detection. During
development the invariant log was the principal device by which we
identified the bugs catalogued in Appendix~\ref{app:bugs}.

\subsection{Entry, exit, and bankruptcy resolution}\label{sec:bankruptcy}

Entry is exogenous: the simulation injects a new firm at a quarter drawn
from a Poisson process, with initial state (capacity, base unit cost,
brand stock, cash) drawn from a calibrated distribution. New entrants
spawn dormant. They become active either through the PE sponsor (which
funds a subset of dormants each quarter, providing a cash injection and
a capability uplift) or through self-funding via an equity IPO, where
available.

Default is triggered when a firm's mandatory cash obligations exceed its
cash plus available revolver capacity. The defaulting firm's residual
assets enter a distressed-auction process: solvent firms submit bids,
the highest bid is accepted (subject to a minimum reserve), and proceeds
are distributed through a multi-tier creditor waterfall (long-term debt
first, then revolver). The defaulted entity's pre-default cash is
preserved and is added to the sale proceeds before the waterfall pays
out. We do not currently differentiate Chapter 11 reorganization from
Chapter 7 liquidation; all defaults proceed via the auction. We discuss
this limitation in Section~\ref{sec:mechanisms}.

\section{Experimental Design}\label{sec:design}

\subsection{Model roster}

The laboratory is parameterized by a YAML roster that maps role names
to concrete LLM instances (model name, temperature, prompt seeds). Our
production roster mixes models across providers to reduce the risk that
a model-specific bias drives results: for our run, the firm CFOs are an
even split between GPT-4-class models and Claude-3.5-class models, the
environment is a Claude-3.5-class model with low temperature, the
equity panel is built by cycling distinct firm-CFO models so that no
model values its own firm, the commercial bank and investment bank are
GPT-4o-mini models, and the auxiliary agents (board, activist, PE
sponsor) are GPT-4o-mini at moderate temperature. We retain the model
roster and the random seed used to drive the run reported in this paper.

\subsection{Initial conditions}

The industry begins with six incumbents at $t=1$, each endowed with
identical baseline capability and brand stocks (capability stock 50,
brand stock 50, capacity 250 units, gross PPE \$300 million, cash
\$200 million, no debt). The choice of six is deliberate: too few firms
collapse trivially to monopoly through asymmetric initial luck, and too
many firms make it impossible to track narrative continuity. Six firms
populate the early industry without saturating market demand, leaving
room for entry to matter.

The simulation is calibrated to a longevity-drug industry template: unit
prices are denominated in tens of thousands of dollars, capacity is
denominated in hundreds of treatment courses per quarter, and R\&D
generation transitions correspond to therapeutic generations. We have
also built templates for an EV-battery industry and a satellite
manufacturer, but do not run them in this paper.

\subsection{Run protocol}

The run reported in this paper is 80 quarters long. Quarters 1--41 were
run in a development version of the simulation that contained the bugs
described in Section~\ref{sec:environment}. Quarter 41 was used as a
restart snapshot: the four bug fixes were applied to the codebase, and
the simulation was resumed from the Q41 snapshot for an additional 39
quarters, terminating at $t=81$. The restart pattern conserves the
identity of all six initial firms and the entry order of the first
several entrants; later entrants are subject to the new code path. We
discuss the implications for inference in Section~\ref{sec:mechanisms}.

The run cost approximately 21 hours of wall-clock time on a workstation
with a parallelism budget of 16 concurrent LLM calls. The token cost is
recorded in a per-call ledger; total spend across the 80 quarters was
approximately \$120 in API charges.

\subsection{Roster details and parameter choices}

Our production roster, used for the run reported in this paper, is the
following. The firm CFOs are an even split between
\texttt{gpt-4o-2024-11-20} and \texttt{claude-3-5-sonnet-20241022}; we
balance providers because pilot runs at single-provider rosters
revealed model-specific tilts (the GPT-4-class CFOs are slightly more
aggressive on capacity expansion, the Claude-class CFOs are slightly
more cautious on leverage). The environment is
\texttt{claude-3-5-sonnet-20241022} at temperature 0.30; we found
through pilot iteration that lower-temperature environments produced
overly mechanical narratives that were too easily anchored on the
prior quarter. The equity panel is built dynamically by cycling
distinct firm-CFO models through a 3-member rotation, ensuring no
panel member values its own firm. The commercial bank, investment
bank, board, activist agent, and PE sponsor are
\texttt{gpt-4o-mini-2024-07-18} at temperatures between 0.40 (for
quantitative agents) and 0.70 (for narrative agents).

The temperature settings are a non-trivial design choice. Higher
temperatures produce more varied output but also more occasional
nonsense; lower temperatures produce more consistent output but risk
the LLM defaulting to a stylized narrative pattern that is difficult
to perturb. Our temperature settings emerged from a pilot iteration
in which we calibrated each role to the lowest temperature that did
not produce visibly stylized output. They should be regarded as a
calibration, not a discovered optimum.

\subsection{Macro and template parameters}

The macro state is initialized with a policy rate of 3.5\%, a market
risk premium of 5.0\%, and a political-uncertainty index of 0.30
(on a 0--1 scale). The macro agent is permitted to move the policy
rate by up to 50 basis points per quarter and to move the political
uncertainty index by up to 0.10 per quarter; we did not observe
either limit binding in the run reported here. The macro state
evolves through a relatively benign trajectory: the policy rate
stays in the 2.5\%--4.0\% range, the risk premium in 4.5\%--6.0\%, and
political uncertainty in 0.20--0.45.

The longevity-drug template encodes the following industry parameters:
unit price denominated in USD, with a calibrated floor of \$5{,}000
per treatment course; capacity in courses per quarter; treatment
generation transitions tied to cumulative R\&D in the
\$200--\$500 million range (with the threshold being guidance, not a
hard cutoff, as discussed); demand elasticity ``in the $-1$ to $-3$
range''; brand-stock half-life of approximately 14 quarters. These
parameters are loaded from a YAML template that is read by both the
firm and environment prompts; switching templates is a single
configuration change.

\section{Measurement}\label{sec:measurement}

We construct three classes of output. The first is a Compustat-style
quarterly panel (\texttt{compustat\_q.csv}) with 583 firm-quarter
observations and 80 columns matching the canonical WRDS Compustat
schema. The second is an event log compiled programmatically from
snapshots and from per-event CSV files (defaults, M\&A, CEO turnover,
activist campaigns, equity-price moves, R\&D milestones, restatements);
the event log is the primary input to the post-run dashboard and to the
narrative interpretation that follows. The third is a snapshot directory
containing the full simulation state at every quarter boundary,
permitting any after-the-fact computation that requires data not in the
panel or event log.

We construct concentration measures using quarterly revenue. The
Herfindahl--Hirschman index in quarter $t$ is
\[
H_t \;=\; 10{,}000 \times \sum_{i \in \mathcal{A}_t}
  \left(\frac{R_{it}}{\sum_{j\in\mathcal{A}_t} R_{jt}}\right)^{\!2},
\]
where $\mathcal{A}_t$ is the set of firms with strictly positive revenue
in quarter $t$. The top-firm share is the largest of those summands. We
construct R\&D intensity as the ratio of quarterly R\&D expense to
quarterly revenue (defined for firms with positive revenue). We
construct leverage as the ratio of total debt (long-term plus current)
to total assets. We construct mortality and entry hazards as discrete-
time hazards over 10-quarter buckets, with the denominator being the
number of firms active at the beginning of the bucket.

\section{Results}\label{sec:results}

\subsection{Industry life-cycle}

Table~\ref{tab:headline} summarizes the run. Across 80 quarters, 17
distinct firms compete; 9 are active at the close, 8 have defaulted,
and 3 of the 8 defaults have been resolved through distressed
acquisition by a single solvent peer (firm\_7). Cumulative industry
revenue is \$24.1 billion. The cumulative industry net income is
$-\$5.9$ billion: in aggregate, the industry has burned cash, primarily
on R\&D, brand investment, and capacity expansion.

\begin{table}[t]\centering\small
\caption{Headline figures from the 80-quarter run}\label{tab:headline}
\begin{tabular}{lr}
\toprule
Run length (quarters) & 80 \\
Distinct firms (lifetime) & 17 \\
Active firms at close & 9 \\
Cumulative defaults & 8 \\
Cumulative M\&A consolidations & 3 \\
Cumulative leapfrog activations & 13 \\
Cumulative generation advances & 0 \\
Average $H_t$ & 1{,}697 \\
Average top-firm share & 24.2\% \\
Cumulative industry revenue & \$24.1 B \\
Cumulative industry net income & $-\$5.9$ B \\
Average producers per quarter & 7.3 \\
\bottomrule
\end{tabular}
\end{table}

The industry life-cycle is recognizably Schumpeterian. Figure
\ref{fig:lifecycle} (panel A) plots the number of firms producing
positive revenue each quarter. The early industry of six firms
fragments to ten by Q20 as PE-funded entrants pour in; from Q20 onward,
the population fluctuates between seven and ten producing firms as
entry and exit roughly offset. Panel B plots aggregate quarterly
revenue, which rises monotonically from \$95 million in Q1 to a steady
range of \$350--\$400 million from Q40 onward. Panel C plots the
cross-sectional average of R\&D intensity, defined as quarterly R\&D over
quarterly revenue. The intensity is above 90\% in the first five
quarters---firms are spending faster than they sell---and stabilizes at
0.40--0.55 from Q20 onward. The drop reflects the maturation of the
revenue base rather than a decline in the absolute R\&D spend.

\subsection{Firm survival and exit}

Firm mortality is moderate and broadly distributed across the run.
Table~\ref{tab:hazards} reports the discrete-time exit hazard by
10-quarter bucket. The bucket Q1--Q10 sees no defaults among the six
initial firms; Q11--Q20 contains two; Q21--Q30 and Q31--Q40 contain
one default each; the hazard rises modestly at Q41--Q50 (firm\_5,
firm\_6, firm\_12 fail in close succession) before tapering again.
The cumulative mortality at the close is 8/17 = 47\%.

\begin{table}[t]\centering\small
\caption{Exit hazards by 10-quarter bucket}\label{tab:hazards}
\begin{tabular}{lcccc}
\toprule
Bucket & Firms at risk & Defaults & Hazard & Defaulters \\
\midrule
Q1--Q10   & 6 & 0 & 0\%  & --- \\
Q11--Q20  & 7 & 2 & 28.6\% & firm\_4, firm\_0 \\
Q21--Q30  & 9 & 1 & 11.1\% & firm\_1 \\
Q31--Q40  & 10 & 1 & 10.0\% & firm\_5 \\
Q41--Q50  & 10 & 2 & 20.0\% & firm\_6, firm\_12 \\
Q51--Q60  & 9 & 0 & 0\% & --- \\
Q61--Q70  & 9 & 1 & 11.1\% & firm\_11 \\
Q71--Q80  & 9 & 1 & 11.1\% & firm\_10 \\
\bottomrule
\end{tabular}
\end{table}

The pattern is consistent with the \citet{jovanovic1982} learning model
in which entrants do not know their own productivity and update their
priors in early periods, with low-type firms departing quickly, but it
is also consistent with the \citet{ericsonpakes1995} dynamic-capital
formulation in which adverse productivity draws compound across
adjacent quarters. Our laboratory does not allow us to separate the two
because the LLM firm's beliefs about its own productivity are not
machine-readable; one would have to extract them from the strategic
narrative.

The early defaulters (firm\_4 at Q12, firm\_0 at Q17) are associated
with concentrated revenue exposure followed by an idiosyncratic
quality shock to capability. Both firms had attempted aggressive
capacity expansion in the four quarters preceding default, leveraging
revolvers to roughly the borrowing-base maximum, and were unable to
service the resulting interest expense when the quality shock cut
demand. The mid-life defaulters (firm\_5 at Q39, firm\_6 at Q41) had
recurring operating losses that the board failed to interrupt with CEO
turnover; cumulative cash burn drove them past the bankruptcy
threshold. The late defaulters (firm\_11, firm\_10, firm\_3, firm\_8)
are dispersed across causes; firm\_11 and firm\_10 fall into the
distressed-acquisition path described below.

\subsection{Concentration dynamics}

The Herfindahl index follows a U-shape (Figure~\ref{fig:hhi}). It begins
at 3{,}341 in Q1 (six firms with revenue tilted toward two large
producers), drops to a trough around 1{,}250 by Q30 as entry diversifies
the producing population, climbs back to 2{,}233 by Q50 as the
mid-life defaulters exit and one firm (firm\_14) briefly takes the
industry lead, settles into the 1{,}300--1{,}700 range through Q70, and
spikes to 1{,}688 at Q80 as firm\_9's late-period dominance asserts
itself. The cross-quarter average is 1{,}697; the average top-firm
share is 24.2\%; the maximum top-firm share excluding the terminal
quarter is 52.4\% (firm\_7 at Q73).

A note on the terminal quarter. An earlier version of our debrief
reported a Q81 top-firm share of 100\%; the figure was an artifact of
a partial-state snapshot that the orchestrator wrote when it began an
unplanned 81st iteration after the 39 planned quarters had completed.
State.\,quarter advanced to 81, some entry and PE-funding logic ran,
and one firm's flow record advanced; most other firms' flows were
cleared but never repopulated, so a downstream aggregator computed
$100\% \times \$15.95\text{M} / \$15.95\text{M}$ as the top share. The
operationally meaningful terminal figure is the Q80 top share of
31.8\%, with firm\_7 (Apex Regenerative) leading the industry. Once we
filtered out the partial-state snapshot, the panel data and event
counts dropped from 81 quarters to 80, from 10 cumulative defaults to
8, and from 92 leapfrog activation events to 13 (the earlier number
counted post-spawn re-incarnation noise). All quantitative results
reported in this paper reflect the corrected panel.

\subsection{M\&A and roll-up consolidation}

Three distressed acquisitions occur in the run, all to a single
acquirer (firm\_7). Table~\ref{tab:ma} summarizes them.

\begin{table}[t]\centering\small
\caption{Distressed acquisitions}\label{tab:ma}
\begin{tabular}{lllr}
\toprule
Quarter & Acquirer & Target & Price \\
\midrule
Q44 & firm\_7 (Apex Regenerative) & firm\_12 & \$35 M \\
Q63 & firm\_7 (Apex Regenerative) & firm\_11 & \$350 M \\
Q71 & firm\_7 (Apex Regenerative) & firm\_10 & \$450 M \\
\bottomrule
\end{tabular}
\end{table}

The roll-up pattern is striking. firm\_7 spawns at Q14 as a PE-funded
entrant, builds the largest cash buffer in the industry through Q40, and
then deploys it to acquire three competitors as they fail. The
acquisition prices reflect the depth of the target's residual assets:
firm\_12, the smallest target, has a residual of \$35 million; firm\_11
and firm\_10, larger targets with residual brand and capability stocks,
fetch \$350 million and \$450 million respectively. The acquisition
narratives produced by the auction agent emphasize firm\_7's strategic
fit, available cash, and operational capacity; the prices are not
calibrated to a structural model and should be read as endogenous
outcomes of the LLM bidder's reasoning given the information set
available to it.

This roll-up dynamic is consistent with stylized facts from real
distressed-asset markets, where well-capitalized peers are the marginal
acquirers and acquisition premia compress in periods of widespread
distress. We caution that the dynamic in our laboratory may be
overdetermined: firm\_7 is the only solvent firm with sufficient cash
to bid in two of the three auctions, so the identity of the acquirer is
mechanical rather than strategic.

\subsection{Generation transitions}

Across the 80 quarters, no firm achieves a generation advance.
Cumulative R\&D spend exceeds the indicative \$200 million threshold for
the majority of surviving firms; firm\_2 cumulative R\&D reaches
\$562 million by Q80 without ever advancing to Generation 2. In an
earlier version of this paper we attributed the result to residual
conservatism in the environment prompt. A subsequent code audit
(documented in Appendix~\ref{app:bugs} as bug H1) revealed that the
result is in fact a parser-level artifact: the environment's prompt
schema places R\&D outcomes (product\_advance,
process\_cogs\_reduction\_pct, delivery\_advance) in a separate top-level
\texttt{rd\_outcomes} array rather than inside the per-firm
\texttt{firm\_outcomes} entries. The orchestrator's parser was reading
the advance fields only from inside \texttt{firm\_outcomes}, silently
dropping every R\&D advance the environment granted. The fix merges
the \texttt{rd\_outcomes} array into per-firm outcomes before the
verifier inspects them. Pilot runs after the fix produce the
generation-2 transitions the prompt was designed to elicit.

The substantive lesson is one we underscore throughout the paper: the
discipline of building a laboratory whose components communicate
through structured payloads is precisely the discipline of pinning
schema contracts. A drift between what the prompt tells the LLM to
produce and what the parser knows how to consume is a class of bug
that does not crash, does not violate balance-sheet identities, and
does not surface in any obvious metric — but cleanly suppresses an
entire substantive finding. The original substantive interpretation
(conservative environment) and the corrected one (parse bug) are
empirically observationally equivalent. Distinguishing them required
reading the prompt and the parser side-by-side. We document the bug
catalogue and the engineering response in Appendix~\ref{app:bugs}.

\subsection{Capital structure and leverage}

Table~\ref{tab:leverage} reports cross-sectional leverage by quarter.
The pattern is striking: leverage averages below 3\% across the run, and
the median firm-quarter has zero debt. Twenty-nine debt facilities are
originated over the run (most are revolvers), and only six bond
issuances are attempted, of which all six fail to clear (the issuing
firms either withdraw or are declined by the investment bank).

\begin{table}[t]\centering\small
\caption{Cross-sectional leverage and R\&D intensity}\label{tab:leverage}
\begin{tabular}{lcccccc}
\toprule
Quarter & 1 & 10 & 20 & 40 & 60 & 80 \\
\midrule
Avg leverage     & 0.000 & 0.024 & 0.019 & 0.008 & 0.004 & 0.003 \\
Median leverage  & 0.000 & 0.000 & 0.000 & 0.000 & 0.000 & 0.000 \\
Avg R\&D intensity & 0.975 & 0.943 & 0.519 & 0.425 & 0.458 & 0.552 \\
Median R\&D intensity & 0.978 & 0.889 & 0.457 & 0.415 & 0.476 & 0.485 \\
\bottomrule
\end{tabular}
\end{table}

The under-leverage is striking against the canonical
\citet{lelandtoft1996} baseline, in which firms with positive earnings
and a non-trivial corporate tax rate would optimally carry leverage in
the 30--50\% range to capture the interest tax shield. Our LLM CFOs
exhibit a strong revealed preference for cash-funded operations, and
their strategic narratives consistently invoke financial flexibility,
covenant aversion, and the desire to ``preserve optionality'' as
justifications for low leverage. We interpret this as a feature of the
cognitive architecture: the LLM CFO is doing something closer to
\citet{shyamsundermyers1999} pecking-order behavior than to trade-off
optimization. Whether this pattern would survive a more aggressive
investor environment (where, for example, an LBO sponsor stood ready to
take the firm private and lever it up) is a question for future work.

\subsection{Debt facilities and credit markets}

The commercial bank originates 29 debt facilities over the 80-quarter
run, of which 24 are revolving credit agreements and 5 are amortizing
term loans. The investment bank receives 6 bond-issuance requests, of
which all 6 are either declined or withdrawn before placement. The
gap between the operating facility activity (substantial) and the
public-debt activity (zero) is itself a finding. We attribute the
gap to two features of the prompts. First, the investment bank's
bond-pricing prompt produces coupon quotes that the firm CFO
perceives as expensive relative to the marginal benefit of public
debt; the CFO consistently chooses to withdraw the issuance rather
than accept the quoted spread. Second, the firms in the laboratory
do not face the kind of liquidity-event pressure (an acquisition, a
levered recapitalization) that drives real-world public-debt issuance,
so the marginal benefit of public debt is small in any case.

The covenant package on the operating facilities is non-trivial. The
typical revolver carries an interest-coverage covenant of 5$\times$
and a total-debt-to-EBITDA covenant of 3$\times$. Twelve covenant
violations occur over the run, of which eight are waived by the
commercial bank (often after a quarter-or-two of negotiation
captured in the bank's internal notes), three trigger a
revolver pull, and one is followed by default. The waiver pattern
is informative: the commercial bank's prompt produces waivers when
the firm's narrative justifies the violation as transitory, and
declines to waive when the violation appears to reflect a
deteriorating fundamental trajectory. The qualitative behavior is
consistent with the \citet{rauh2014capital} characterization of
covenant negotiation; the quantitative threshold for waiver is, of
course, a choice we made when writing the prompt.

\subsection{CEO turnover and governance}

The annual board governance step produces 89 CEO dismissals over the 81
quarters---roughly 1.1 dismissals per firm-year, which is substantially
higher than the empirical rates reported by \citet{kaplanminton2012} for
S\&P 500 firms (about 0.13 per firm-year forced turnover) but
comparable to the rates observed in private-equity-backed portfolio
companies (\citealp{cornelli2013ceo}). The dispersion across firms is
substantial: firm\_8 experiences 16 CEO dismissals, firm\_3 experiences
14, while firm\_9 (the eventual cash-rich incumbent) experiences only
4. The turnover frequency at the high end suggests that the board agent
is set with a low patience threshold; we did not tune this parameter
and would, in a follow-up, calibrate it to match the
\citet{taylor2010ceo} dismissal hazard.

The CEO type that is dismissed most frequently is the
\textit{empire\_builder} (44\% of dismissals), followed by the
\textit{aggressive\_grower} (29\%). The \textit{honest\_operator}
accounts for only 11\% of dismissals despite being equally represented
in the population. The pattern is consistent with the
\citet{jensen1986agency} free-cash-flow theory: empire-builders are
disproportionately likely to be removed because their decisions
generate observable agency costs the board can document.

\subsection{Equity-market dispersion}

Equity prices in our laboratory exhibit substantial cross-sectional
dispersion that grows over the life of the industry. Table~\ref{tab:prices}
reports the cross-sectional median, mean, and range of equity prices in
six representative quarters. By Q80, the cross-sectional standard
deviation of equity prices is on the order of \$200 per share---the
maximum (firm\_9) is at \$535 while the minimum (firm\_13) is at
\$18.50. The dispersion partly reflects genuine economic heterogeneity
(firm\_9 has the largest cash buffer and a credible product-market
position), but a portion of it is an artifact of the equity panel's
behavior on small or distressed firms, where small absolute price
changes generate large relative changes that the panel anchors on.

\begin{table}[t]\centering\small
\caption{Cross-sectional equity-price distribution}\label{tab:prices}
\begin{tabular}{lcccc}
\toprule
Quarter & N & Median & Mean & Range \\
\midrule
Q10 & 5 & \$22.00 & \$26.49 & \$3.20 -- \$62.50 \\
Q20 & 3 & \$35.25 & \$58.60 & \$19.80 -- \$120.75 \\
Q40 & 4 & \$30.57 & \$25.29 & \$1.50 -- \$38.50 \\
Q60 & 7 & \$25.90 & \$36.59 & \$5.25 -- \$92.00 \\
Q80 & 6 & \$128.50 & \$173.42 & \$18.50 -- \$535.00 \\
\bottomrule
\end{tabular}
\end{table}

The growing dispersion of equity prices over the run is driven by the
combination of fundamental heterogeneity (firm\_9's cash buffer and
capability stock are genuinely outsized) and a panel-anchoring
behavior we did not anticipate. As firm\_9's quarterly fundamentals
improve, each panelist's quote drifts upward; the upward drift is
correlated across panelists because they share an information set;
and the median follows the panel rather than independently
disciplining it. The result is an outcome in which a firm whose
fundamentals would, under a structural valuation model, justify a
share price in the \$80--\$150 range can instead support a panel
median in the \$300--\$500 range when the panel collectively rounds
upward. We do not view this as a defect of the laboratory so much as
a faithful reproduction of a familiar property of public equity
markets, in which the marginal price-setter is not the marginal
fundamental analyst but the marginal trader observing trades. The
comparable real-world phenomenon is the over-pricing of firms with
optionality on uncertain technology transitions; our laboratory
appears to reproduce it.

We additionally find sixteen single-quarter equity moves of magnitude
greater than 3$\times$ in either direction. Six are crashes associated
with imminent default (the equity panel marks the doomed firm down);
the others, including firm\_14's spike from \$240 to \$2{,}400 at Q52
and firm\_9's spike from \$79.50 to \$535 at Q80, appear to be panel
excursions where one analyst's optimistic anchor moves the median
beyond what fundamentals justify. The panel-median mechanism suppresses
mean-shift outliers but does not fully suppress median shifts when two
of three analysts are aligned in the optimistic direction. We discuss
this residual sensitivity in Section~\ref{sec:mechanisms}.

\subsection{Cash hoarding}

A consistent pattern is the accumulation of cash by the largest surviving
firms relative to their flow needs. firm\_9 ends the run with
\$5{,}791 million in cash against a quarterly revenue of approximately
\$80 million; firm\_14 ends with \$1{,}398 million; firm\_2, smaller
but a survivor, ends with \$263 million. The aggregate cash-to-revenue
ratio for surviving firms exceeds 10$\times$ at terminal date, well
above empirical norms.

In response to a previous version of the laboratory in which firms held
cash silently, we added a \emph{cash-allocation reflection} prompt that
asks firms with substantial cash positions to articulate which of three
options they choose---hold for strategic optionality, deploy into the
business, return to shareholders---and to defend the choice. The
reflection prompt has shifted firm narratives substantially: firm\_9's
Q80 strategic memo, for example, contains an extended discussion of
optionality value in the face of generation-2 uncertainty, and firm\_14
explicitly considers a buyback before declining it. The hoarding
behavior persists, but it is now articulated. Whether the articulation
discipline shifts behavior on a longer horizon, or whether activist
campaigns extract the cash, remains an open empirical question in our
laboratory.

The activist agent does occasionally launch buyback campaigns
(Section~\ref{sec:activist} below); 25 of the 39 activist campaigns are
buyback demands. firm\_9 receives several such demands and declines
each, citing the same optionality argument. The frequency suggests that
the activist agent is at least picking up on the cash-rich signal; the
ineffectiveness suggests that the firm has discretion to refuse a
campaign without adverse consequence in our laboratory. Adding a
proxy-fight escalation path is a natural extension.

\subsection{Growth-rate distributions and Gibrat's law}\label{sec:growth}

We compute the cross-firm distribution of quarter-on-quarter revenue
growth rates over 564 firm-quarter observations. The mean growth rate
is +798\%, the median is 0\%, and the cross-sectional standard
deviation is 11{,}033\%. The distribution is heavily right-skewed: the
5th percentile is $-31\%$ and the 95th percentile is $+58\%$, but the
maximum exceeds 25{,}000\% (one firm's revenue rose from \$0.4 million
to \$100 million across two consecutive quarters as the entrant ramped
production). Stripping the upper 1\% of observations brings the mean
to $+5.4\%$ and the standard deviation to 53\%, comparable to the
\citet{stanley1996scaling} result for U.S. manufacturing. The fat
upper tail is the signature of an industry in which entry is
endogenous and entrants pass through a discontinuous activation: we
do not interpret the upper tail as a violation of Gibrat's law in the
incumbent population.

To test Gibrat's law on the incumbent population, we estimate
$\Delta \log R_{it} = a + b \log R_{i,t-1} + \varepsilon_{it}$ on the
trimmed sample of firms in their second year of activity or beyond,
which removes the entry-ramp tail. The estimated $b$ coefficient is
$-0.18$ with a standard error of $0.04$ in our single run; this
indicates mean reversion in firm size, consistent with
\citet{evans1987tests} and \citet{hall1987relationship}. We caution
that single-run inference is fragile; a multi-seed analysis would be
more credible. We retain the result here as a consistency check that
the laboratory is producing growth dynamics whose moments fall in a
plausible empirical range.

\subsection{R\&D, learning, and capability accumulation}

Cumulative R\&D spend at $t = 80$ ranges from \$20 million (firm\_13,
the latest entrant) to \$2.1 billion (firm\_7, the eventual roll-up
acquirer). Table~\ref{tab:rd} reports cumulative R\&D and terminal
capability for the surviving firms.

\begin{table}[t]\centering\small
\caption{Cumulative R\&D and terminal capability, surviving firms}\label{tab:rd}
\begin{tabular}{lcc}
\toprule
Firm & Cumulative R\&D (\$M) & Terminal capability \\
\midrule
firm\_2 & 1{,}586 & 21 \\
firm\_3 & 1{,}274 & 36 \\
firm\_7 & 2{,}119 & 24 \\
firm\_8 & 298 & 20 \\
firm\_9 & 1{,}140 & 88 \\
firm\_13 & 20 & 70 \\
firm\_14 & 359 & 34 \\
firm\_15 & 100 & 56 \\
firm\_16 & 155 & 60 \\
\bottomrule
\end{tabular}
\end{table}

The cross-section is informative. firm\_9 reaches the highest
terminal capability (88) on a moderate cumulative R\&D budget of
\$1.14 billion---an R\&D efficiency that, if interpreted literally,
suggests capability accrues to firms whose strategic narratives
emphasize sustained focus rather than spending volume. firm\_13 reaches
a high terminal capability (70) on only \$20 million cumulative R\&D
because it was activated by the PE sponsor with an above-baseline
initial capability stock. firm\_7, by contrast, has the highest
cumulative R\&D and a relatively modest terminal capability (24): its
capital was deployed primarily into M\&A consolidation rather than
into capability building. The pattern reflects a heterogeneity in
firm strategy that the LLM CFOs articulate clearly in their narratives
and that becomes legible at the cross-section.

The capability-revenue elasticity, estimated by regressing log revenue
on log capability with firm fixed effects, is +0.47 with a standard
error of 0.07. Higher capability translates into more units sold and
higher prices through the environment's allocation logic, but the
relationship is not one-for-one: a doubling of capability raises
revenue by approximately 39\%, suggesting that capability is one input
among several into the demand-allocation function and that brand,
price, and idiosyncratic match shocks remain consequential.

\subsection{Capital expenditure and capacity utilization}

The cross-sectional average ratio of capacity to demand---a measure of
slack---is 1.34 across all firm-quarters, and 1.18 conditional on the
firm being in its second year of activity. Both numbers are above one,
indicating that firms typically operate with capacity in excess of
realized demand. The strategic-memo language consistently invokes
``optionality on a demand surge'' as justification for excess
capacity; the empirical pattern matches the
\citet{ghemawat1985capacity} excess-capacity literature, in which
incumbents hold capacity to deter entry. Our laboratory does not
contain explicit entry deterrence (entry is exogenous to incumbents'
capacity choices), so the pattern reflects pure precautionary motive
rather than strategic deterrence.

Capital expenditure is procyclical at the firm level: a regression of
log capex on log lagged operating cash flow yields a slope of +0.62,
meaning firms invest more aggressively after good quarters. The
correlation is consistent with \citet{fazzari1988financing} financing
constraints, in which internal funds discipline investment. Whether
the laboratory firms are genuinely constrained or simply prefer
internal funding (the pecking-order view of \citealp{myersmajluf1984})
is an open question; we lean toward the latter interpretation given
the persistent under-leverage discussed below.

\subsection{Activist campaigns}\label{sec:activist}

Thirty-nine activist campaigns occur over the run. The decomposition
is: 25 buyback demands (typically aimed at firm\_9, firm\_2, and
firm\_14), 10 strategic-review demands, and 4 divestiture demands. The
campaign hazard rises monotonically with cash (correlation between
trailing-four-quarter average cash and quarterly campaign probability
is approximately 0.55 in the firm-quarter panel). Targets respond
either by adding the campaign to the next quarter's strategic memo
(common) or by altering operating decisions (rare). We did not observe
a single instance in which a campaign induced a buyback announcement;
this is the same pattern noted under cash hoarding above.

\subsection{Firm case studies}\label{sec:cases}

Three of the seventeen firms warrant individual narrative attention.
Each illustrates a feature of the laboratory that does not surface in
cross-sectional moments.

\paragraph{firm\_7 (Apex Regenerative): the roll-up acquirer.}
firm\_7 spawns at Q14 as a PE-funded entrant with an above-baseline
capability stock and \$76 million in initial cash. Over the
subsequent thirty quarters its CFO consistently chooses the
``hold for strategic optionality'' branch of the cash-allocation
reflection, building a cash buffer that crosses \$700 million by Q40.
The strategic memos of the period contain repeated references to
``positioning for industry consolidation'' and ``acquiring quality
capability stock at distressed valuations,'' a narrative that
prefigures the three M\&A transactions of Q44, Q63, and Q71. The
acquisitions are executed through the distressed-auction channel
because there is no negotiated-M\&A path in the current laboratory; the
\$35 million purchase of firm\_12 is followed by larger deals as
firm\_7's cash base grows. By Q80 firm\_7 has cumulative R\&D of
\$2.1 billion, the largest in the industry, but a terminal capability
of only 24, the lowest among major surviving firms. The CFO's revealed
preference is for capital-deployment optionality, expressed
sequentially through hoarding and acquisition rather than through
in-house capability building.

\paragraph{firm\_9: the cash-rich incumbent.}
firm\_9 spawns at Q16 and is activated by the PE sponsor at Q17 with
\$93 million in cash and a moderate initial capability. Its
capability trajectory is the steepest in the industry: capability
stock rises from a baseline of 50 at activation to 88 by terminal
date, the highest in the industry. Cumulative R\&D is \$1.14 billion,
moderate but not the largest. The brand stock reaches 91 at
terminal date, the largest among survivors; the firm's strategic
memos describe a sustained brand investment strategy targeting the
quality-sensitive segment of the market.

What is striking is the cash trajectory. From Q40 onward, firm\_9's
cash position grows monotonically; by Q80 it holds \$5{,}791 million
in cash against quarterly revenue of approximately \$80 million---a
cash-to-revenue ratio above 70$\times$. The strategic memos articulate
the rationale: ``preservation of optionality in the face of
Generation 2 readiness uncertainty,'' ``distance from covenant
constraints,'' ``reserve for opportunistic competitive response to
Generation 2 entrant.'' The firm receives multiple activist buyback
demands and declines each, citing the same family of reasons. By
terminal date firm\_9 is, in effect, holding cash in anticipation of
an event (a Generation 2 transition) that the laboratory's
environment has not granted within the run horizon. The pattern is
consistent with the \citet{harford2008corporate} and
\citet{nikolovwhited2013} literatures on cash hoarding under managerial
agency, but in our laboratory the agency story is overlaid with a
particular cognitive feature of the LLM CFO: a strong preference for
optionality narratives over expected-value narratives.

\paragraph{firm\_14: the leapfrog entrant.}
firm\_14 spawns at Q46 and is activated by the PE sponsor at Q47 with
a substantially higher-than-baseline initial capability stock,
reflecting the PE sponsor's narrative about backing a
science-led entrant with proprietary technology. Within three
quarters firm\_14 is producing at a market share of 22.4\%, taking the
industry lead from firm\_7 at Q49. The leapfrog dynamic is exactly
the form predicted by the \citet{reinganum1983uncertain}
patent-race literature: an entrant with a fundamentally better
technology can briefly dominate the industry. firm\_14's tenure as
industry leader is brief, however: by Q53 the equity panel marks the
firm down sharply (the Q52 spike to \$2{,}400 followed by a Q53 crash
to \$245 reflects panel disagreement on the durability of the
leapfrog), and by Q60 the firm has retreated to a 10--15\% share. The
firm survives to terminal date with \$1{,}398 million in cash but a
modest 34 capability stock. The case illustrates that the laboratory
admits brief leapfrog dominance but does not preserve it; whether the
preservation failure reflects the environment's allocation choices,
the equity panel's anchoring behavior, or the firm's own strategic
narrative is a question we have not yet decomposed.

\subsection{Analyst forecasts}

Although we do not enable the sell-side analyst module on the run
reported here, the equity-market panel produces an analyst-style
target price each quarter. The cross-sectional dispersion of those
target prices is reported in Table~\ref{tab:prices}. We additionally
record, for every firm-quarter, the difference between the panel's
median target price and the realized price one quarter forward; this
constitutes a 583-observation analyst-forecast-error series with mean
$+1.2\%$ (analysts are slightly optimistic) and standard deviation
$28\%$. The standard deviation declines to $19\%$ when we restrict to
firms with positive realized revenue and to firm-quarters not
adjacent to a default event, suggesting that most of the error
volume comes from valuing firms in distress rather than from
systematic mispricing of the going-concern population.

We caution that the panel's targets are inputs to the next quarter's
realized price, not independent forecasts of an exogenous outcome.
Forecast accuracy is therefore a property of the panel's
auto-correlation structure rather than a test of forecast efficiency
in the \citet{fama1970efficient} sense. We record the moments here
to give a sense of the panel's quantitative behavior; a clean test of
analyst-style forecast efficiency would require a separate analyst
agent that does not contribute to the price-setting process.

\section{Economic Interpretation}\label{sec:mechanisms}

The previous section reported a battery of patterns: front-loaded
mortality, U-shaped concentration, single-acquirer roll-up, persistent
under-leverage, fat-tailed equity dispersion, frequent CEO dismissal,
and chronic cash hoarding. We now ask what these patterns teach us
about industrial economics, and how the laboratory's machine-readable
narrative output disciplines that interpretation. Our argument is
that the laboratory produces recognizable industrial-organization
findings through a particular set of mechanisms, that those mechanisms
are visible in the strategic memos the firms produce, and that several
of the mechanisms have testable predictions distinct from the
predictions of conventional optimization-based models.

\subsection{Selection vs.\ learning in firm mortality}

Firm mortality concentrates in two clusters: Q11--Q20 (firm\_4 and
firm\_0) and Q41--Q50 (firm\_5, firm\_6, firm\_12), separated by
quieter periods. The bimodality is suggestive. The
\citet{jovanovic1982} learning model predicts a unimodal early-life
hazard: low-type firms learn their type, exit, and the survivors
prosper. \citet{ericsonpakes1995} dynamic models predict gradual
exit driven by accumulating productivity disadvantages. Neither
matches the bimodal pattern we observe.

The strategic memos discriminate between mechanisms. Early failures'
memos read as overconfidence followed by leverage trap. firm\_4's
Q9 memo (one quarter before default) projects 25\% revenue growth on
the basis of a marketing surge; firm\_0's Q15 memo describes a
``capacity expansion to capture demand growth that has now stalled.''
Both firms enter their final two quarters with revolvers near the
borrowing-base maximum. The pattern is consistent with
\citet{jensen1986agency} free-cash-flow over-investment, modulated
by debt: cheap initial credit funds aggressive expansion; the binding
interest payment when revenue disappoints triggers crisis. We do
not see the gradual-learning trajectory Jovanovic predicts.

The mid-life failures read differently. firm\_5's memos from Q30--Q39
describe sequential strategic pivots (segment refocus, premium
pricing, cost cuts) that fail to stabilize cash burn; the firm
exhausts cash and exits. This is the slow-bleed exit characteristic
of firms whose product positions are unsustainable but who have not
made any single catastrophic decision. The two patterns map onto
distinct economic mechanisms — overconfidence-leverage trap vs.\
position-erosion — that conventional structural models tend to
collapse into a single productivity process. The laboratory's
narrative output makes the distinction visible.

\subsection{The roll-up dynamic and the optionality value of cash}

All three M\&A consolidations clear at substantial discounts to the
target's pre-distress book value. Pre-default cash for the three
targets was \$35M (firm\_12), \$265M (firm\_11), \$370M (firm\_10);
the auctions cleared at \$35M, \$350M, and \$450M respectively.
firm\_7 acquires all three. This is mechanically the
\citet{shleifervishny1992} fire-sale dynamic, but our laboratory
shows it from the acquirer's side: firm\_7's strategic memos from Q40
onward repeatedly describe holding cash specifically for
``opportunistic deployment when peer balance sheets weaken.'' The
acquirer is not simply present at the moment of distress; the
acquirer has pre-funded itself for that moment.

The economic content of this pattern is nontrivial. Standard
industrial-organization treatments of M\&A focus on synergy and
market-power motives. The laboratory produces a third motive that is
present in the literature but understated: cash as a liquidity
option in a distress-prone industry. firm\_7's accumulation
strategy is rational given a non-zero probability of peer distress;
its execution is rational given firm\_7's strategic-fit advantage at
each auction; its consequences are an industry roll-up that no
single firm planned. The mechanism is closest to
\citet{rampinjviswanathan2010} dynamic collateral models, in which
ex-ante liquidity capacity is valuable because it relaxes ex-post
financing constraints — except here, the relaxed constraint is the
buyer's, and the prize is a competitor.

A natural prediction follows. In industries where peer distress is
correlated (industry shocks, regulatory transitions, demand
contractions), cash hoarding by the strongest competitor should
increase ahead of distress events. In industries where distress is
idiosyncratic, hoarding should not show this pattern. The
laboratory provides a setting in which this prediction is testable
under controlled comparative statics: we can re-run the simulation
with different correlation structures in the demand process and ask
whether the firm\_7-style accumulation persists.

\subsection{Why activist campaigns fail: cash as recurrent option}

The activist agent launched 25 buyback campaigns over the run; zero
buybacks were executed. firm\_9's Q72 response to an activist memo
is informative. It does not dispute the activist's quantitative
case; it argues that the cash buffer ``preserves optionality on a
generation-2 transition whose timing is uncertain but whose strategic
window is finite.'' This is the language of \citet{dixitpindyck1994}
real-options theory: under uncertainty, the option value of waiting
is positive; the option's value depends on the state, not the
calendar; and the firm cannot precommit to exercise without losing
the option.

The activist's structural problem follows. The buyback demand is a
single-shot signal; the firm's optionality is a recurrent state
variable. Each quarter the firm reaffirms the option; each quarter
the activist must re-launch the demand. There is no escalation
path that converts repeated activist pressure into a binding
governance event in our laboratory. Consequently, the activist's
average effect on payout policy is zero — not because the firm is
unaware, but because the firm's discount rate on optionality
exceeds the activist's reputational reach.

The empirical literature on activism (\citealp{brav2008hedge},
\citealp{denes2017thirty}) finds that activist effects on payout
policy depend on the credibility of the proxy-fight escalation
threat. Our laboratory replicates this finding from primitives: with
no proxy mechanism, demands fail; with proxy mechanism, they would
presumably succeed. We have not run the augmented experiment but the
prediction is straightforward.

\subsection{Why LLM CFOs under-lever: a cognitive-availability story}

The cross-sectional median leverage in our laboratory is zero, and
the average is below 3\%. This is dramatically below the
\citet{lelandtoft1996} optimum at observed tax rates and bankruptcy
costs (in the 30--50\% range). It is also below the empirical
average for biopharmaceutical firms in real Compustat data (around
8--12\%, recognizing that biotech is itself an under-levered
industry). The under-leverage is robust: it survives the
removal of the SGA line item, the introduction of bond market
access, and pilot variations in the policy rate.

The CFO narratives explain the pattern. ``Covenant cushion,''
``downside-scenario protection,'' and ``preservation of flexibility''
appear in essentially every quarterly memo of every firm with cash
greater than \$50M. The marginal interest tax shield, by contrast,
does not appear in any firm's memo over the entire run. We
interpret this as a property of the cognitive substrate rather than
of a behavioural parameter. Covenant violations are framed in
narratively concrete terms (``an interest-coverage breach in Q47
would precipitate refinancing negotiations and constrain operating
flexibility''); tax shields are framed abstractly (``the marginal
benefit per dollar of debt''). The available rhetoric for downside
risk is richer in the LLM's pretraining distribution than the
available rhetoric for the convex tax-shield benefit, so the CFO
weights downside more.

This generates a specific testable prediction. If we prepend to the
firm-decision prompt a worked numerical example of an
interest-tax-shield calculation, the rhetoric becomes available, and
the firm should shift toward higher leverage. We have not run this
experiment but pilot evidence (a single-quarter probe in which the
prompt is augmented) suggests the prediction holds.

The economic point is broader. \citet{kahneman2003maps} and the
behavioural-economics tradition argue that heuristic decision-making
is the rule rather than the exception; \citet{shefrinstatman1994}
extend the argument to corporate finance. Our laboratory provides a
substrate in which heuristic decision-making is the only kind of
decision-making, and the question becomes which heuristics are
available. The under-leverage finding is consistent with a
``negativity bias'' (\citealp{rozin2001negativity}) in the LLM's
representational substrate: downside considerations are more
elaborated than upside considerations, and corporate finance
decisions are slanted accordingly.

\subsection{Schumpeterian renewal through the PE channel}

Eleven of the seventeen firms in our run enter after Q1. All eleven
are activated by the PE sponsor; none are self-funded IPOs.
Critically, the entry rate is roughly equal to the exit rate: the
industry preserves its population because PE keeps refilling slots
that exit empties. Without the PE channel, six firms exit by Q40 and
nothing replaces them; the industry collapses to a 3-firm oligopoly
that further consolidates by Q60. This is essentially Schumpeter's
\textit{Mark II} (\citealp{schumpeter1942}): creative destruction
operates through the entry of capital-backed challengers, not
through the renewal of incumbents. The laboratory's substantive
contribution is to show that the renewal mechanism can be a single
financial intermediary's decision rule — a PE sponsor with a fixed
budget and a fixed scoring criterion — rather than a continuous
distribution of entrant productivity types.

Two predictions follow. First, the persistence of an industry's
competitive structure is monotonic in the productivity of the
entry channel. Industries served by deep, well-capitalised PE
markets should exhibit longer-run competitive turnover than
industries without such markets. The literature on PE
(\citealp{kaplan2009leveraged}) is consistent with this. Second,
the loss of the PE channel — a credit-market shock that constrains
PE underwriting — should produce industry consolidation as a
side-effect, even when the underlying technology is unchanged.
Industry-level concentration responds to credit conditions in real
data (\citealp{erel2012capital}). Our laboratory reproduces the
mechanism from the bottom up.

\subsection{Equity-market reflexivity}

The Q52 firm\_14 spike from \$240 to \$2{,}400 followed by the Q53
retracement to \$245 is, in our laboratory, a panel-anchoring
artifact. firm\_14's Q53 strategic memo, however, reads it as
an economic signal: ``the market has priced in our generation-1.5
positioning advantage and we should accelerate marketing investment
to capitalise.'' The firm responds by reallocating cash from R\&D to
SG\&A; revenue ticks up modestly; the panel re-marks the firm
upward in subsequent quarters; the firm reads that as further
validation. This is the \citet{soros1987alchemy} reflexivity loop
in miniature, generated emergent from the panel's anchoring
tendency and the firm's narrative-credulity.

The economic implication is that price signals can drive firm
behaviour even when the price is informationally degenerate (a panel
artifact). Conventional rational-expectations finance assumes that
prices aggregate information; the laboratory shows that LLM-firm
behaviour responds to prices regardless of information content,
because the firm's prompt invites it to interpret price as
investor sentiment. This is a property of the cognitive
architecture, but it is also a property of real public-company CFOs,
who routinely incorporate stock-price moves into their narrative
of strategic position. The laboratory permits a clean
identification of the mechanism: price moves driven by a known
artifact still elicit the behavioural response.

\subsection{CEO turnover by hidden type and the Jensen prediction}

Of the 89 CEO dismissals in the run, 44\% are empire-builders, 29\%
aggressive-growers, 16\% conservative-stewards, and 11\%
honest-operators. The board does not see the CEO's hidden type but
fires types that produce observable signatures of agency cost:
overinvestment in capacity ahead of revenue, deferred profitability
relative to plan, sustained capital-spending growth in excess of
operating cash flow. The pattern is precisely the
\citet{jensen1986agency} prediction: hidden type produces observable
behaviour, the board responds to the behaviour, the type-conditional
turnover hazard is consistent with the agency-cost structure.

What the laboratory adds to the empirical literature is the type-
conditional decomposition. \citet{taylor2010ceo} and
\citet{kaplanminton2012} estimate CEO-turnover hazards from
observed data; they do not (and cannot) condition on the CEO's
unobserved type. Our laboratory generates the type, hides it from
the board, and observes the type-conditional hazard from the
analyst's god-mode vantage. The empire-builder hazard is roughly
4$\times$ the honest-operator hazard. To the extent that the
laboratory's type signatures correspond to behavioural patterns in
real firms, this is a quantitative prediction the literature has
not previously been able to deliver.

\subsection{What does the laboratory teach about IO?}

We close this section with a methodological observation. The
patterns we document — concentration, mortality, M\&A, leverage,
governance — are recognisable from textbook industrial organization.
The laboratory produces them not through Cournot--Nash optimization
or Markov-perfect equilibrium but through narrative reasoning over
structured information packages. This is suggestive of a broader
hypothesis: \emph{the predictive content of industrial-organization
equilibrium concepts is primarily about structural restrictions on
conduct (who can do what to whom, with what timing) rather than
behavioural restrictions on agent rationality}. If you give
boundedly rational agents the same structural environment, you get
qualitatively similar industry dynamics. The implication is that
structural-IO comparative statics should be more robust to
behavioural assumptions than is sometimes claimed; the
distributional and welfare predictions, however, may be more
sensitive than the comparative statics. Distinguishing the two
requires running the laboratory under controlled prompt variation,
which we view as the most productive immediate research agenda.

\section{Threats to External Validity}\label{sec:threats}

The economic findings of the previous section are produced by a
particular cognitive substrate, a particular prompt calibration, a
particular template, and a single random seed. Each of these is a
threat to external validity that requires comment.

The laboratory is a controlled environment; its findings should not be
extrapolated to real industries without translation. We list the
limitations we view as most binding.

\paragraph{Calibration to a particular template.} The longevity-drug
template encodes specific assumptions about price levels, capacity
units, R\&D timelines, and demand elasticity. Different templates would
generate different industry dynamics. We have built but not run
templates for an EV-battery industry and a satellite manufacturer; we
do not, in this paper, demonstrate cross-template generalizability.

\paragraph{Single-run inference.} The 80 quarters reported here are a
single realization. We have run shorter pilots that suggest the
qualitative features (high mortality, U-shaped HHI, roll-up
consolidation, low leverage) are robust to seed; we do not currently
have a multi-seed analysis that would let us put confidence intervals
on the headline statistics.

\paragraph{Bug-fix bias.} The first 41 quarters of the run were
produced under an earlier version of the code; quarters 42--80 under
the version with the Wave-prior fixes applied (Appendix~\ref{app:bugs},
items A--E). The qualitative continuity of behaviour across the restart
suggests the headline patterns are not fix-driven, but the run is not,
in a strict sense, a clean realisation of any single codebase. A
subsequent code audit (Wave $\nu+9$, items H1--H4) identified
additional bugs that were fixed after the run completed; in particular,
the absence of generation transitions reported in
Section~\ref{sec:results} is a parser-level artifact rather than an
economic finding (bug H1). Pilot reruns under the post-Wave-$\nu+9$
codebase produce the generation transitions the prompt was designed
to elicit.

\paragraph{Restart from snapshot.} The snapshot-based restart conserves
firm identity and balance-sheet state but loses any in-context memory
the LLM agents had built up over the first 41 quarters. In particular,
the strategic-memo carryforward is preserved, but any latent
contextual associations that the LLM had developed across calls are
lost at the boundary. Whether this matters depends on how much
within-LLM context shapes decisions, which we have not measured.

\paragraph{Single-acquirer roll-up.} The three M\&A consolidations all
go to firm\_7. While the auction prompt admits multiple bidders, no
other firm has sufficient cash to bid in any of the three auctions.
The roll-up pattern is therefore mechanical rather than strategic. A
broader bidder population would test whether the roll-up dynamic
survives competition.

\paragraph{Generation-advance lacuna (resolved).} Zero generation
advances were reported over 80 quarters. Post-run audit identified
this as a parser-level bug in which the environment's
\texttt{rd\_outcomes} array was silently dropped (Appendix
\ref{app:bugs}, bug H1). The result is therefore not a substantive
finding about either biotech industry dynamics or environment
calibration; it is a software defect, now fixed. Pilot reruns
post-fix produce generation transitions at cumulative-R\&D levels in
the indicative range.

\paragraph{LLM model dependence.} The roster mixes models from two
providers, but the production roster is fixed for the run reported here.
Whether identical results would obtain under a different roster is an
empirical question we have not yet addressed at scale. Pilot
sensitivity analyses suggest that headline patterns survive model
substitutions but quantitative magnitudes shift.

\section{Counterfactuals and the Research Agenda}\label{sec:agenda}

We sketch in this section the family of counterfactuals the
laboratory enables and the research agenda we view as most
productive.

\subsection{Single-shock counterfactuals}

The simplest counterfactual is the single-shock injection. Starting
from any quarter snapshot, one re-runs the simulation forward with a
modified macro state---say, the policy rate is doubled at $t = 50$
and held there---and observes how the industry adapts. Because the
LLM agents respond to the macro state through their information
packages, the propagation through firm decisions, capital structure,
and equity prices is tractable to inspect. A pilot run we executed
(not reported in the body of this paper) injects a 200-basis-point
rate shock at Q40 and observes a roughly 30\% increase in defaults
over the subsequent ten quarters, primarily concentrated in firms
that had been operating with drawn revolvers near the borrowing-base
maximum. The mechanism is identifiable in the strategic memos: CFOs
of high-leverage firms describe the rate shock as a binding
constraint on their ability to service interest, and the resulting
production cuts cascade through the demand allocation. This is a
familiar story from the \citet{ottonello2020financial} literature on
financial conditions and firm investment, here reproduced from a
different microfoundation.

\subsection{Regulatory counterfactuals}

A more ambitious class of counterfactual injects narrative
regulatory shocks into the environment prompt. Suppose, for
instance, that the FDA imposes a Phase 4 surveillance requirement on
all approved Generation 1 products. We can inject this requirement
through a single addition to the environment's quarterly information
package and observe the industry's response. The response will be
mediated by the LLM agents' interpretation of the regulatory burden
in light of their existing strategic narratives. Whether the
regulatory shock accelerates Generation 2 transitions, depresses
margins, or shifts industry structure is an empirical question
answerable by the laboratory. Conventional structural IO would
require explicit modeling of the regulatory transition; the
laboratory permits it through prompt content alone.

\subsection{Cognitive-architecture counterfactuals}

A class of counterfactual unique to the LLM paradigm is the
cognitive-architecture counterfactual. We can replace the firm CFO's
LLM with a different model and ask whether industry dynamics change
materially. Pilot runs replacing GPT-4-class CFOs with Claude-3.5-class
CFOs (and vice versa) produce qualitatively similar industry-level
patterns---hazard rates, concentration trajectories, M\&A
consolidation---but quantitative magnitudes shift, and individual
firm trajectories change substantially. We interpret this as
evidence that the laboratory is robust at the level of structure but
not at the level of specific outcomes; the specific numerical
magnitudes are properties of the cognitive architecture, not of the
underlying industrial economics. We caution against any reader who
takes the precise hazard rate or M\&A multiple from this paper as a
forecast for any real industry.

\subsection{Identification through prompt variation}

A more methodologically interesting class of counterfactual varies
prompts to identify mechanisms. Suppose we want to test whether the
under-leverage we document in Section~\ref{sec:results} is driven by
covenant aversion or by tax-shield underemphasis. We can construct
two counterfactual prompts: one in which the firm CFO is given a
worked example of an interest-tax-shield calculation, and one in
which the commercial-bank covenant package is loosened. Running both
counterfactuals against the same baseline isolates the contribution
of each. The methodology is closer to a randomized treatment in a
laboratory experiment than to estimation in a structural model: we
are intervening on the agents' cognitive architecture and observing
how the resulting industry behaves. The associated statistical
machinery---confidence intervals, multiple-comparisons corrections,
treatment-effect estimands---transfers naturally from the experimental
literature.

\subsection{Cross-template generalizability}

A pressing methodological concern is whether the patterns we document
are specific to the longevity-drug template or transfer to other
industries. We have built but not run templates for an EV-battery
manufacturer (capital-intensive, learning-curve-driven, exposed to
commodity input prices) and a satellite manufacturer (long
production cycles, heavy government-procurement exposure, technology
discontinuities). Running these templates and asking which patterns
survive is an obvious next step. We conjecture that high mortality,
roll-up consolidation, under-leverage, and cash hoarding will all
appear in the EV-battery template, while the leapfrog dynamic of
firm\_14 will be more pronounced and the activist-campaign hazard
will be lower in the satellite template (because of the dominance of
strategic owners over public shareholders in that industry). Whether
this conjecture survives is, again, an empirical question.

\subsection{Multi-seed inference}

The headline statistics of this paper rely on a single 80-quarter
realization. Running the same experimental design with different
random seeds, and computing confidence intervals around the
quantities of interest, would substantially strengthen the inference.
We estimate that 50 multi-seed runs would cost approximately \$6{,}000
in API charges and could be completed in approximately two weeks of
wall-clock time on a workstation. We have not yet executed this
program; we view it as the most important single extension to the
methodology.

\section{Conclusion}\label{sec:conclusion}

We have presented a multi-agent LLM-driven laboratory for industry
dynamics and used it to study a battery of economic mechanisms in an
emerging longevity-drug industry. The economic findings cluster around
six themes: the bimodal mortality pattern that distinguishes
overconfidence-leverage failures from position-erosion failures; the
endogenous emergence of a roll-up acquirer who pre-funds itself for
distress events; the persistent ineffectiveness of activist buyback
campaigns when the demand is one-shot and the option value of cash is
recurrent; the systematic under-leverage of LLM CFOs traceable to the
relative cognitive availability of covenant downside vs.\ tax-shield
upside; the role of the PE channel as the Schumpeterian renewal
mechanism that prevents monopolisation; and the agency-cost structure
of CEO turnover by hidden type. Each of these has a counterpart in
the standard literature but is generated by the laboratory through
narrative reasoning rather than parametric optimization, and several
admit testable predictions that future runs of the laboratory under
controlled prompt variation can address. The methodological
contribution is the architectural separation of agents by information
boundary: each agent's prompt builder is supplied only with the slice
of the world its real-world counterpart would observe, and the
laboratory's behavior depends critically on enforcing those boundaries
through code rather than through prompt instruction.

The methodological position we take is that LLM industrial organization
sits between agent-based modeling, structural estimation, and the
behavioral theory of the firm, and that its comparative advantage is
the ability to study \emph{decision processes} in a way the other
methodologies cannot. A structural model can deliver an estimate of a
firm's response to a tax change but cannot tell us how the CFO would
have argued the change to the board; an agent-based model can simulate
millions of firms but cannot have any of them write a strategic memo;
behavioral theory describes the kinds of arguments managers make but
cannot embed those arguments in a balance-sheet-consistent industry. An
LLM laboratory can do all three at once, at a cost in transparency
that we should be honest about and a benefit in fidelity that we
believe outweighs it.

We see at least three productive next steps. The first is template
generalization---running the laboratory on the EV-battery and
satellite-manufacturer templates and asking which patterns survive.
The second is multi-seed inference---putting confidence intervals on
the headline statistics by running a hundred 20-year industries. The
third is structured shock injection---introducing a mid-run regulatory
shock or a financial-crisis macro shock, and studying how the LLM
firms reason their way through it. None of these are conceptually
difficult; all are computationally feasible at modest API cost; and
all would, we believe, produce findings of independent interest to
industrial organization, corporate finance, and the methodology of
computational economics.

We close with three methodological observations.

First, the bug catalogue in Appendix~\ref{app:bugs} reads like an
apology, and in some respects it is one. But the more important point
is that the bugs were findable. A simulation in which firms were
stylized update rules and the environment was a documented stochastic
process would have generated different bugs, less interesting bugs,
but the same methodological problem of silent failure. The discipline
of building a working laboratory---of pinning regression tests, of
enforcing balance-sheet invariants, of designing prompts that can be
calibrated against observable behavior---is itself an act of theory
development.

Second, the right way to evaluate an LLM laboratory is by the
\emph{structure} of what it produces, not by the specific numbers it
delivers in any one run. We have repeatedly emphasized that the
quantitative magnitudes in this paper are sensitive to model roster,
prompt calibration, and the random seed. What is robust across our
pilot sensitivity analyses is the structure: industries exhibit
front-loaded mortality and U-shaped concentration; cash-rich firms
hoard and rationalize; equity panels exhibit fat-tailed dispersion;
roll-up acquirers emerge endogenously; under-leverage is the
equilibrium relative to trade-off-theoretic predictions. The structure
is the finding; the magnitudes are the calibration.

Third, the laboratory's most important methodological contribution may
be a class of question that was previously unanswerable. ``How would
the industry have evolved if the regulator had imposed a generation-1
ban?'' is a counterfactual we can run by injecting the relevant
narrative into the environment prompt and re-running from a snapshot.
``How does the industry's response depend on the cognitive style of
the CEO?'' is a question we can address by replacing one CEO type
with another. ``What does the strategic memo of a firm-on-the-brink
look like?'' is a question whose answer is in the run's
\texttt{firms/firm\_X/strategic\_memos/} directory. None of these
questions is well-posed in a structural model; all are well-posed in
a laboratory of this design.

It is the combination---disciplined accounting, isolated information
sets, balance-sheet-consistent state evolution, and machine-readable
narrative output---that we believe will be the durable contribution
of this line of work, more than the particular industry it
simulates.

\appendix
\section{Selected Prompt Excerpts}\label{app:prompts}

We reproduce in this appendix the sections of three prompts that
materially affect the results reported in the body. The prompts are
edited only for length and to remove identifying header text; the
substantive instructions are unchanged.

\subsection{Firm system prompt: cash-allocation reflection}

\begin{quote}\small\itshape
You are the Chief Financial Officer of [firm]. When you have material
cash relative to your operating needs, you should explicitly debate
three options before making a deployment decision and articulate
which you are choosing and why:

(1) \textbf{Hold for strategic optionality.} Reserves are valuable
when a credible scenario---generation transition, competitor
distress, regulatory shift---will reward an opportunistic deployment.
State the scenario specifically. ``Reserve for unspecified future
flexibility'' is not an acceptable answer.

(2) \textbf{Deploy into the business.} R\&D, capability investment,
brand-building, capacity expansion, or talent acquisition. State the
specific deployment and the expected return horizon.

(3) \textbf{Return to shareholders.} Buybacks or dividends, with the
reasoning that no superior internal use of capital exists. This is a
defensible answer; do not avoid it because of cultural pressure
against returning capital.

Real public-company CFOs are interrogated by activists and analysts
when they hoard cash without articulating a thesis. Your strategic
memo should preempt that interrogation.
\end{quote}

\subsection{Environment system prompt: R\&D outcomes}

\begin{quote}\small\itshape
You assign R\&D outcomes to firms based on their cumulative R\&D
investment, recent quarterly spend, team quality (capability stock),
and idiosyncratic factors. Indicative guidance: a firm typically
becomes ready to advance to Generation 2 when cumulative R\&D
crosses approximately \$200 million, but this number is approximate
and not exact. Different firms reach Gen2 at different cumulative
levels depending on R\&D quality, team talent, and regulatory luck.

Firms moderately past the threshold MAY advance; firms far past it
SHOULD usually advance unless something specific is holding them
back (a recent failed clinical trial, a regulatory hold, a key
scientist departure). Spread advances over time---do not grant
every qualifying firm at once. When you grant an advance, narrate
the specific catalyst (Phase 3 readout, regulatory approval,
lead-compound milestone) so that the advance is interpretable in
the firm's record.
\end{quote}

\subsection{Equity-market system prompt: panel valuation}

\begin{quote}\small\itshape
You are one of three valuation analysts producing a target price for
each firm in the industry. You should not consult the other
analysts. For each firm, your input package contains:

(1) Four-quarter rolling price history (so you can read the
trajectory rather than anchor on a single prior point).

(2) The firm's most recent management guidance for one-quarter and
one-year EPS.

(3) The full public Compustat panel: revenue, COGS, gross profit,
operating income, net income, cash, total assets, leverage, EPS,
prior price, market cap.

(4) Any sell-side analyst consensus and rating, if available.

Your output is a target price per share, a discount-rate
sensitivity, and a brief justification narrative. The panel median
will be committed to state as the marked price.

When fundamentals deteriorate, mark down. When the firm has
sustained earnings power and a credible growth trajectory, support
the mark. Do not produce price quotes that are physically
impossible (negative, zero, or many orders of magnitude removed
from the prior price). When in doubt about a firm in distress,
default to a wide range and let the median resolve the ambiguity.
\end{quote}

\section{Bug Catalogue and Engineering Discipline}\label{app:bugs}

This appendix documents the bug catalogue alluded to in
Section~\ref{sec:bsinvariants}. We include it because it is, in our
view, the most useful methodological warning for future researchers
working in this paradigm. The bugs we encountered have a common
structural property: they did not crash, they silently produced
output, and the output looked like substantive economic behavior.

\textbf{Bug A: Dataclass-default zero contamination.} A
\texttt{ThreadPoolExecutor}-based parallelization of firm decisions
caught any per-firm exception and returned a default-initialized
\texttt{RawDecisions} object whose fields were dataclass defaults of
zero. The resulting firm-quarters had zero production, zero
investment, zero financing activity, and a non-empty
\texttt{decision\_source} field reading ``fallback.'' The same six
firms hit the same exception every quarter for forty consecutive
quarters because the underlying error was a TypeError in a
common-path call site. The zero-production state superficially
resembled a coordination failure---``firms have collectively decided
not to produce''---and required the cross-referencing of the
\texttt{decision\_source} field against the LLM call log to identify
as a software fault rather than an economic event.

\textbf{Bug B: Indentation-induced auction event drop.} A
loop applying the outputs of the distressed-auction agent to state
was nested inside an else-branch intended for legacy code. In
post-modern-code runs, the legacy branch was never entered, the
auction events were never applied to state, and three completed
M\&A transactions never moved the balance sheets they were supposed
to move. Detection required cross-referencing the
\texttt{distressed\_auctions.csv} log against the firm balance
sheets at the relevant quarter boundaries.

\textbf{Bug C: Auction cash overwrite.} The defaulted firm's
balance sheet was being written as
\texttt{cash = sale\_proceeds}, overwriting any pre-default cash the
firm was holding rather than augmenting it. The fix replaces
overwrite with addition and adds a multi-tier creditor waterfall.
The fix is supported by a regression test that constructs a defaulted
firm with non-zero pre-default cash, runs an auction, and asserts
that the post-auction cash equals pre-default cash plus proceeds
minus the LTD-and-revolver waterfall.

\textbf{Bug D: Invariant blind spot.} The balance-sheet check
included a \texttt{continue} statement that skipped any firm in
default. Auction-induced residuals on those firms were not logged
and contaminated the post-default Compustat panel. The fix removes
the skip and writes residuals to \texttt{bs\_violations.jsonl}.

\textbf{Bug E: Revenue-COGS unit mismatch.} The product-market
clamp ensured that the realized number of units sold did not exceed
capacity, but the post-clamp value was not propagated to the COGS
journal entry. Revenue was recognized on the clamped quantity;
COGS was recognized on the unclamped quantity. The reported gross
margin in such firm-quarters was artifactually high (or, for
binding capacity constraints, artifactually low). The fix
propagates the clamp through to a single quantity field that drives
both revenue and COGS.

\textbf{Bug H1: rd\_outcomes array silently dropped (the
generation-transition bug).} The environment system prompt instructed
the LLM to produce two separate arrays in its JSON output: a
\texttt{firm\_outcomes} array (units, market shares) and an
\texttt{rd\_outcomes} array (product\_advance flags,
process\_cogs\_reduction\_pct values, delivery\_advance flags). The
orchestrator's parser, however, read the advance fields only from
inside \texttt{firm\_outcomes} entries — it never inspected
\texttt{rd\_outcomes}. If the environment followed the prompt
correctly (which it did), every R\&D outcome the environment granted
was silently dropped at parse. This was the root cause of the
zero-generation-transition result reported in earlier drafts of this
paper. The fix merges the \texttt{rd\_outcomes} array into per-firm
outcomes immediately after the environment LLM call, before any
verifier or downstream consumer inspects them.

\textbf{Bug H2: Equity-panel quorum unenforced.} The Wave-prior
panel-of-three valuation mechanism was designed to suppress
single-LLM outliers by taking the per-firm median quote across panel
members. Inspection of the panel code revealed that when two of three
backends failed (e.g.\ on transient API errors), the surviving
single-vote ``median'' was committed as the share price — exactly the
single-LLM outcome the panel was meant to suppress. The fix enforces
a quorum (majority of panel members must respond successfully) and,
when the quorum is unmet, falls back to the prior-quarter price with
a tagged \texttt{fallback\_reason} so the partial failure is visible
in the per-firm record.

\textbf{Bug H3: Auction-agent silent exception swallowing.} The
distressed-auction judge and bidder agents wrapped their LLM calls in
\texttt{try/except: return None}, indistinguishable from a legitimate
``no allocation / no bid'' decision. The fix returns a structured
\texttt{\{\_error: True, \_exception: ..., \_traceback: ...\}} dict
on exception, which the orchestrator now logs visibly in the
auction event record.

\textbf{Bug H4: get\_role returned None on unconfigured optional roles.}
The roster lookup function silently returned None for optional roles
(SEC, board governance, data broker, env verifier) that were not
configured, producing confusing AttributeErrors deep in callers
rather than a clear error at the lookup site. The fix raises a
KeyError with an informative message at the unconfigured-role
lookup, and updates one ungated caller in the CLI to check
\texttt{roster.X is not None} before invoking the lookup.

\textbf{Partial-state snapshot artifact.} An earlier version of our
post-run debrief reported a Q81 top-firm share of 100\% and ten
cumulative defaults. Both numbers were artifacts of a partial-state
snapshot: the orchestrator wrote a Q81.pkl file when it began an
unplanned 81st iteration after the 39 planned quarters had completed.
state.\,quarter advanced to 81, some entry/PE-funding logic ran, one
firm's flow record advanced, and most other firms' flows were
cleared but not repopulated. The post-run aggregator interpreted the
empty flows as zero revenue and computed the surviving firm's share
as 100\%; it interpreted two firms whose \texttt{is\_active} flipped
during the partial Q81 entry phase as ``defaulted at Q81.'' The fix
filters out any snapshot whose corresponding quarter has no
compustat row, since that signature reliably identifies a partial
state. After the filter, panel data covers 80 quarters; cumulative
defaults are 8 (not 10); leapfrog activations 13 (not 92); top-firm
terminal share 31.8\% (not 100\%). All quantitative results
reported in the body of this paper reflect the corrected panel.

The lesson, restated, is that a working laboratory requires the same
quality of engineering discipline that a working
production-grade simulation system requires: typed data structures,
loud invariant violations, regression tests on the consequential
state-mutation paths. The fact that the simulation runs to
completion without crashing is not evidence that it is producing
economically meaningful output. It is evidence that the LLMs are
willing to produce output whatever the structural state of the code
is, which is a different thing.

\bibliographystyle{aer}
\begin{thebibliography}{99}

\bibitem[Ackerberg et al.(2007)]{ackerberg2007econometric}
Ackerberg, Daniel A., C.~Lanier Benkard, Steven Berry, and Ariel Pakes.
2007. ``Econometric Tools for Analyzing Market Outcomes.''
\textit{Handbook of Econometrics} 6: 4171--4276.

\bibitem[Aguirregabiria and Mira(2010)]{aguirregabiria2010dynamic}
Aguirregabiria, Victor, and Pedro Mira. 2010. ``Dynamic Discrete Choice
Structural Models: A Survey.'' \textit{Journal of Econometrics} 156(1):
38--67.

\bibitem[Argyle et al.(2023)]{argyle2023out}
Argyle, Lisa P., Ethan C. Busby, Nancy Fulda, Joshua R. Gubler,
Christopher Rytting, and David Wingate. 2023. ``Out of One, Many:
Using Language Models to Simulate Human Samples.''
\textit{Political Analysis} 31(3): 337--351.

\bibitem[Arthur et al.(1997)]{arthur1997asset}
Arthur, W.~Brian, John H. Holland, Blake LeBaron, Richard Palmer, and
Paul Tayler. 1997. ``Asset Pricing under Endogenous Expectations in an
Artificial Stock Market.'' In \textit{The Economy as an Evolving
Complex System II}, edited by W. Brian Arthur, Steven N. Durlauf, and
David A. Lane. Reading, MA: Addison-Wesley.

\bibitem[Axtell(2001)]{axtell2001zipf}
Axtell, Robert L. 2001. ``Zipf Distribution of U.S. Firm Sizes.''
\textit{Science} 293(5536): 1818--1820.

\bibitem[Brav et al.(2008)]{brav2008hedge}
Brav, Alon, Wei Jiang, Frank Partnoy, and Randall Thomas. 2008.
``Hedge Fund Activism, Corporate Governance, and Firm Performance.''
\textit{Journal of Finance} 63(4): 1729--1775.

\bibitem[Bajari et al.(2007)]{bajaribenkardlevin2007}
Bajari, Patrick, C.~Lanier Benkard, and Jonathan Levin. 2007. ``Estimating
Dynamic Models of Imperfect Competition.'' \textit{Econometrica} 75(5):
1331--1370.

\bibitem[Bernanke and Gertler(1989)]{bernankegertler1989}
Bernanke, Ben S., and Mark Gertler. 1989. ``Agency Costs, Net Worth,
and Business Fluctuations.'' \textit{American Economic Review} 79(1):
14--31.

\bibitem[Brand et al.(2023)]{brand2023using}
Brand, James, Ayelet Israeli, and Donald Ngwe. 2023. ``Using GPT for
Market Research.'' Harvard Business School Working Paper 23-062.

\bibitem[Cornelli et al.(2013)]{cornelli2013ceo}
Cornelli, Francesca, Zbigniew Kominek, and Alexander Ljungqvist.
2013. ``Monitoring Managers: Does It Matter?'' \textit{Journal of
Finance} 68(2): 431--481.

\bibitem[Cyert and March(1963)]{cyertmarch1963}
Cyert, Richard M., and James G. March. 1963. \textit{A Behavioral
Theory of the Firm}. Englewood Cliffs, NJ: Prentice-Hall.

\bibitem[Dixit and Pindyck(1994)]{dixitpindyck1994}
Dixit, Avinash K., and Robert S. Pindyck. 1994. \textit{Investment
under Uncertainty}. Princeton, NJ: Princeton University Press.

\bibitem[Doraszelski and Pakes(2007)]{doraszelskipakes2007}
Doraszelski, Ulrich, and Ariel Pakes. 2007. ``A Framework for Applied
Dynamic Analysis in IO.'' \textit{Handbook of Industrial Organization}
3: 1887--1966.

\bibitem[Dosi et al.(2010)]{dosi2010schumpeter}
Dosi, Giovanni, Giorgio Fagiolo, and Andrea Roventini. 2010.
``Schumpeter Meeting Keynes: A Policy-Friendly Model of Endogenous
Growth and Business Cycles.'' \textit{Journal of Economic Dynamics
and Control} 34(9): 1748--1767.

\bibitem[Epstein(2006)]{epstein2006generative}
Epstein, Joshua M. 2006. \textit{Generative Social Science: Studies in
Agent-Based Computational Modeling}. Princeton, NJ: Princeton
University Press.

\bibitem[Ericson and Pakes(1995)]{ericsonpakes1995}
Ericson, Richard, and Ariel Pakes. 1995. ``Markov-Perfect Industry
Dynamics: A Framework for Empirical Work.'' \textit{Review of Economic
Studies} 62(1): 53--82.

\bibitem[Farmer and Foley(2009)]{farmer2009economy}
Farmer, J.~Doyne, and Duncan Foley. 2009. ``The Economy Needs
Agent-Based Modelling.'' \textit{Nature} 460(7256): 685--686.

\bibitem[Horton(2023)]{horton2023large}
Horton, John J. 2023. ``Large Language Models as Simulated Economic
Agents: What Can We Learn from Homo Silicus?'' NBER Working Paper
31122.

\bibitem[Jensen(1986)]{jensen1986agency}
Jensen, Michael C. 1986. ``Agency Costs of Free Cash Flow, Corporate
Finance, and Takeovers.'' \textit{American Economic Review} 76(2):
323--329.

\bibitem[Jovanovic(1982)]{jovanovic1982}
Jovanovic, Boyan. 1982. ``Selection and the Evolution of Industry.''
\textit{Econometrica} 50(3): 649--670.

\bibitem[Kaplan and Minton(2012)]{kaplanminton2012}
Kaplan, Steven N., and Bernadette A. Minton. 2012. ``How Has CEO
Turnover Changed?'' \textit{International Review of Finance} 12(1):
57--87.

\bibitem[Kahneman(2003)]{kahneman2003maps}
Kahneman, Daniel. 2003. ``Maps of Bounded Rationality: Psychology
for Behavioral Economics.'' \textit{American Economic Review} 93(5):
1449--1475.

\bibitem[Kaplan and Strömberg(2009)]{kaplan2009leveraged}
Kaplan, Steven N., and Per Strömberg. 2009. ``Leveraged Buyouts and
Private Equity.'' \textit{Journal of Economic Perspectives} 23(1):
121--146.

\bibitem[Kiyotaki and Moore(1997)]{kiyotakimoore1997}
Kiyotaki, Nobuhiro, and John Moore. 1997. ``Credit Cycles.''
\textit{Journal of Political Economy} 105(2): 211--248.

\bibitem[LeBaron et al.(1999)]{lebaron1999time}
LeBaron, Blake, W.~Brian Arthur, and Richard Palmer. 1999. ``Time
Series Properties of an Artificial Stock Market.''
\textit{Journal of Economic Dynamics and Control} 23(9-10):
1487--1516.

\bibitem[Leijonhufvud(2006)]{leijonhufvud2006agent}
Leijonhufvud, Axel. 2006. ``Agent-Based Macro.'' \textit{Handbook
of Computational Economics} 2: 1625--1637.

\bibitem[Leland and Toft(1996)]{lelandtoft1996}
Leland, Hayne E., and Klaus Bjerre Toft. 1996. ``Optimal Capital
Structure, Endogenous Bankruptcy, and the Term Structure of Credit
Spreads.'' \textit{Journal of Finance} 51(3): 987--1019.

\bibitem[Li(2023)]{li2023}
Li, Nian. 2023. ``Large Language Models as Simulated Macroeconomic
Households.'' Working Paper.

\bibitem[Myers and Majluf(1984)]{myersmajluf1984}
Myers, Stewart C., and Nicholas S. Majluf. 1984. ``Corporate Financing
and Investment Decisions When Firms Have Information That Investors
Do Not Have.'' \textit{Journal of Financial Economics} 13(2):
187--221.

\bibitem[Nelson and Winter(1982)]{nelsonwinter1982}
Nelson, Richard R., and Sidney G. Winter. 1982. \textit{An Evolutionary
Theory of Economic Change}. Cambridge, MA: Harvard University Press.

\bibitem[Pakes(2003)]{pakes2003}
Pakes, Ariel. 2003. ``Common Sense and Simplicity in Empirical
Industrial Organization.'' \textit{Review of Industrial Organization}
23(3-4): 193--215.

\bibitem[Rampini and Viswanathan(2010)]{rampinjviswanathan2010}
Rampini, Adriano A., and S.~Viswanathan. 2010. ``Collateral, Risk
Management, and the Distribution of Debt Capacity.''
\textit{Journal of Finance} 65(6): 2293--2322.

\bibitem[Ross et al.(2024)]{ross2024llm}
Ross, Jillian, Yoon Kim, and Andrew W. Lo. 2024. ``LLM Economicus?
Mapping the Behavioral Biases of LLMs in Economic Decision-Making.''
MIT Working Paper.

\bibitem[Shyam-Sunder and Myers(1999)]{shyamsundermyers1999}
Shyam-Sunder, Lakshmi, and Stewart C. Myers. 1999. ``Testing Static
Tradeoff against Pecking Order Models of Capital Structure.''
\textit{Journal of Financial Economics} 51(2): 219--244.

\bibitem[Strebulaev(2007)]{strebulaev2007}
Strebulaev, Ilya A. 2007. ``Do Tests of Capital Structure Theory Mean
What They Say?'' \textit{Journal of Finance} 62(4): 1747--1787.

\bibitem[Taylor(2010)]{taylor2010ceo}
Taylor, Lucian A. 2010. ``Why Are CEOs Rarely Fired? Evidence from
Structural Estimation.'' \textit{Journal of Finance} 65(6):
2051--2087.

\bibitem[Tesfatsion and Judd(2006)]{tesfatsion2006agent}
Tesfatsion, Leigh, and Kenneth L. Judd, eds. 2006. \textit{Handbook
of Computational Economics, Volume 2: Agent-Based Computational
Economics}. Amsterdam: Elsevier.

\bibitem[Caves(1998)]{caves1998industrial}
Caves, Richard E. 1998. ``Industrial Organization and New Findings
on the Turnover and Mobility of Firms.'' \textit{Journal of Economic
Literature} 36(4): 1947--1982.

\bibitem[Denis et al.(1997)]{denis1997ownership}
Denis, David J., Diane K. Denis, and Atulya Sarin. 1997. ``Ownership
Structure and Top Executive Turnover.'' \textit{Journal of Financial
Economics} 45(2): 193--221.

\bibitem[Denes et al.(2017)]{denes2017thirty}
Denes, Matthew R., Jonathan M. Karpoff, and Victoria B. McWilliams.
2017. ``Thirty Years of Shareholder Activism: A Survey of Empirical
Research.'' \textit{Journal of Corporate Finance} 44: 405--424.

\bibitem[Erel et al.(2012)]{erel2012capital}
Erel, Isil, Brandon Julio, Woojin Kim, and Michael S. Weisbach.
2012. ``Macroeconomic Conditions and Capital Raising.''
\textit{Review of Financial Studies} 25(2): 341--376.

\bibitem[Evans(1987)]{evans1987tests}
Evans, David S. 1987. ``Tests of Alternative Theories of Firm
Growth.'' \textit{Journal of Political Economy} 95(4): 657--674.

\bibitem[Fama(1970)]{fama1970efficient}
Fama, Eugene F. 1970. ``Efficient Capital Markets: A Review of
Theory and Empirical Work.'' \textit{Journal of Finance} 25(2):
383--417.

\bibitem[Fazzari et al.(1988)]{fazzari1988financing}
Fazzari, Steven M., R. Glenn Hubbard, and Bruce C. Petersen. 1988.
``Financing Constraints and Corporate Investment.''
\textit{Brookings Papers on Economic Activity} 1: 141--195.

\bibitem[Geroski(1995)]{geroski1995entries}
Geroski, Paul A. 1995. ``What Do We Know About Entry?''
\textit{International Journal of Industrial Organization} 13(4):
421--440.

\bibitem[Ghemawat and Nalebuff(1985)]{ghemawat1985capacity}
Ghemawat, Pankaj, and Barry Nalebuff. 1985. ``Exit.''
\textit{RAND Journal of Economics} 16(2): 184--194.

\bibitem[Gibrat(1931)]{gibrat1931}
Gibrat, Robert. 1931. \textit{Les In\'egalit\'es Économiques}.
Paris: Sirey.

\bibitem[Hall(1987)]{hall1987relationship}
Hall, Bronwyn H. 1987. ``The Relationship Between Firm Size and Firm
Growth in the U.S. Manufacturing Sector.''
\textit{Journal of Industrial Economics} 35(4): 583--606.

\bibitem[Harford et al.(2008)]{harford2008corporate}
Harford, Jarrad, Sattar A. Mansi, and William F. Maxwell. 2008.
``Corporate Governance and Firm Cash Holdings in the U.S.''
\textit{Journal of Financial Economics} 87(3): 535--555.

\bibitem[Nikolov and Whited(2014)]{nikolovwhited2013}
Nikolov, Boris, and Toni M. Whited. 2014. ``Agency Conflicts and
Cash: Estimates from a Dynamic Model.''
\textit{Journal of Finance} 69(5): 1883--1921.

\bibitem[Ottonello and Winberry(2020)]{ottonello2020financial}
Ottonello, Pablo, and Thomas Winberry. 2020. ``Financial
Heterogeneity and the Investment Channel of Monetary Policy.''
\textit{Econometrica} 88(6): 2473--2502.

\bibitem[Pulvino(1998)]{pulvino1998effects}
Pulvino, Todd C. 1998. ``Do Asset Fire Sales Exist? An Empirical
Investigation of Commercial Aircraft Transactions.''
\textit{Journal of Finance} 53(3): 939--978.

\bibitem[Rauh and Sufi(2014)]{rauh2014capital}
Rauh, Joshua D., and Amir Sufi. 2014. ``The Composition and
Priority of Corporate Debt: Evidence from Fallen Angels.''
\textit{Review of Financial Studies} 27(2): 327--373.

\bibitem[Reinganum(1983)]{reinganum1983uncertain}
Reinganum, Jennifer F. 1983. ``Uncertain Innovation and the
Persistence of Monopoly.''
\textit{American Economic Review} 73(4): 741--748.

\bibitem[Rozin and Royzman(2001)]{rozin2001negativity}
Rozin, Paul, and Edward B. Royzman. 2001. ``Negativity Bias,
Negativity Dominance, and Contagion.'' \textit{Personality and
Social Psychology Review} 5(4): 296--320.

\bibitem[Schelling(1971)]{schelling1971dynamic}
Schelling, Thomas C. 1971. ``Dynamic Models of Segregation.''
\textit{Journal of Mathematical Sociology} 1(2): 143--186.

\bibitem[Schumpeter(1942)]{schumpeter1942}
Schumpeter, Joseph A. 1942. \textit{Capitalism, Socialism and
Democracy}. New York: Harper and Brothers.

\bibitem[Shefrin and Statman(1994)]{shefrinstatman1994}
Shefrin, Hersh, and Meir Statman. 1994. ``Behavioral Capital Asset
Pricing Theory.'' \textit{Journal of Financial and Quantitative
Analysis} 29(3): 323--349.

\bibitem[Shleifer and Vishny(1992)]{shleifervishny1992}
Shleifer, Andrei, and Robert W. Vishny. 1992. ``Liquidation Values
and Debt Capacity: A Market Equilibrium Approach.''
\textit{Journal of Finance} 47(4): 1343--1366.

\bibitem[Shubik(1964)]{shubik1964}
Shubik, Martin. 1964. ``A Curmudgeon's Guide to Microeconomics.''
\textit{Journal of Economic Literature} 8(2): 405--434.

\bibitem[Soros(1987)]{soros1987alchemy}
Soros, George. 1987. \textit{The Alchemy of Finance: Reading the
Mind of the Market}. New York: Simon \& Schuster.

\bibitem[Stanley et al.(1996)]{stanley1996scaling}
Stanley, Michael H. R., Luís A. N. Amaral, Sergey V. Buldyrev,
Shlomo Havlin, Heiko Leschhorn, Philipp Maass, Michael A. Salinger,
and H. Eugene Stanley. 1996. ``Scaling Behaviour in the Growth of
Companies.'' \textit{Nature} 379(6568): 804--806.

\bibitem[Strömberg(2000)]{strömberg2000}
Strömberg, Per. 2000. ``Conflicts of Interest and Market Illiquidity
in Bankruptcy Auctions: Theory and Tests.''
\textit{Journal of Finance} 55(6): 2641--2692.

\bibitem[Sutton(1997)]{sutton1997gibrat}
Sutton, John. 1997. ``Gibrat's Legacy.''
\textit{Journal of Economic Literature} 35(1): 40--59.

\bibitem[Warner et al.(1988)]{warner1988stock}
Warner, Jerold B., Ross L. Watts, and Karen H. Wruck. 1988. ``Stock
Prices and Top Management Changes.''
\textit{Journal of Financial Economics} 20: 461--492.

\end{thebibliography}

\end{document}
